% 牌理（はいり）：手牌の速度と価値を最大化させるための理論
% 牌効率（はいこうりつ）：四面子一雀頭を効率よく完成させるための理論

\documentclass[../nankiru_main.tex]{subfiles}

% override main tex render option
\renewcommand{\renderOption}{0}

\begin{document}

\subfilepreamble{役作り・場況対応}

\subsection{役作り}

ここまでの話の趣旨は「最速聴牌」と「最高の聴牌形」。立直、ドラ、裏ドラ、一発などの役によって他の手役を狙わなくても満貫になる効率がいいので、手役を作る技術の利用頻度が低い。しかし、速度が同じでも、役の選択により手牌の価値が変動する。また、局収支として「最速聴牌」が正解ではない場面もある。特に打点が欲しい場合、または真っ直ぐに進めばドラなしリーのみの一途になる場合、手役を重視して、純粋な牌効率を逆らうことになる。

手役の同士を比較する方法と、どの程度まで手役を狙うために牌効率が逆らえるか、この節で説明していく。


\subsubsection*{麻雀は満貫を作るゲーム}

先におきたい念頭はこのような打点効率の話。麻雀って、満貫以下の打点は1飜があればあるほど打点が2倍になる。4飜ならピンフツモでない限り打点が8000ほどあって、次の12000になるまでもう2飜が必要。この打点上昇の幅は飜数に見合わないため、多くの場合は満貫を優先する。そして、3900などもう1飜で3800の差が出る手牌なら、多少受け入れを減らしてもまだ見合う。5200や6400などの場合なら、もう1飜で打点の差はそほどではないため、「ほぼ満貫」と扱ってもいい。

手牌\ref{mj:3900}は \mj{3p} を切れば3900の手牌になってしまう。\mj{45p} を払うと遅くなるが打点が7700になるので、残りの形は両面両面なら見合う\citep[][pp.~98--99]{fkt}。

\tehai{40m 3345p 22334s -66'6z}[8巡目　ドラ \mj{3p}][3900]


%---------------------------------------------
\subsubsection{役牌} \label{subsub:yakuhai}

役牌が重なると仕掛けができるので、「聴牌への急行券」とも思われる。が、役牌は字牌と同じ、横引きができない。孤立の19牌の11枚の受け入れに対して、孤立の役牌の受けは3枚だけ。それに、生牌字牌が重ならないまま持っていくと、危険度が巡目が深くなるほど高まる。役牌と19牌（または28牌）の比較はこの小節で説明していく。

結論からいうと\textbf{諸説ある}。これは（1）オタ風と役牌と19牌と質は違うからと（2）役牌が他家にも利用できるので自分の手牌・場況・点棒状況などの要素に左右されやすいから。どっちから切るかには統一の答えがない。そして、人によって、同じ牌姿と場況でも、手牌への評価が違って、字牌と19牌の比較も違う。人による比較の違いを理解するには、まず19牌と役牌とオタ風の性質の理解が必要。

\begin{itemize}
  \item オタ風
  \begin{itemize}
    \item 一人の相手にしか利用できない $\to$ 安全度が高い
    \item ピンフの頭として働ける $\to$ ピンフとの相性が高い
    \item 縦引きしかできない $\to$ 手牌を早めない
  \end{itemize}
  \item 役牌
  \begin{itemize}
    \item 重なったら役がある $\to$ 手牌の速度を上げる
    \item 誰にも利用できる $\to$ 相手に鳴かれたくない時は早めに切る（ダブ東ダブ南など尚更）
    \item ポンしたら他の部分に頭が求められる $\to$ ヘッドレスやピンフとの相性が悪い
    \item 振り込んだ単騎でない限り相手に役を与える
  \end{itemize}
  \item 19牌
  \begin{itemize}
    \item 横引きがある $\to$ 効率としては優位だが、両面にはならない
    \item 振り込んだ組み合わせが字牌より多い $\to$ 字牌より危険
  \end{itemize}
\end{itemize}

そしていろんな方の重要度順を挙げると、

\begin{itemize}
  \item \citet{senba-yakuhai}
  \begin{itemize}
    \item 手牌が整っている（対子１つ＋好形2ブロック）：19牌＞オタ風＞役牌
    \item さもないと：役牌＞オタ風＞19牌
  \end{itemize}
  \item \citet{tsuchida-yakuhai}
  \begin{itemize}
    \item アガリそうな、アガりたい手：19牌＞オタ風＞役牌
    \item アガリなさそうな手：オタ風＞役牌＞28牌＞19牌
  \end{itemize}
  \item \citet{naga-yakuhai}によるNAGAの切り順
  \begin{itemize}
    \item 明確なタンピン系手牌：オタ風＞役牌
    \item さもないと：役牌＞19牌＞オタ風
  \end{itemize}
  \item \citet{shibukawa-yakuhai}
  \begin{itemize}
    \item 1345の1＞役牌＞オタ風
    \item オタ風はほぼ第一打
  \end{itemize}
  \item \citet{yoteru-yakuhai}
  \begin{itemize}
    \item ドラが1枚：19牌＞役牌＞オタ風
    \item ドラなし：役牌＞19牌＞オタ風
  \end{itemize}
  \item \citet{chigo-yakuhai}
  \begin{itemize}
    \item ヘッドレスは役牌から
  \end{itemize}
\end{itemize}

これらの重要度順に見えるテーマは、

\begin{itemize}
  \item 門前立直に近いか
  \item タンピン系なのか
  \item ドラありかなしか
\end{itemize}

%\textbf{牌効率からいうと、19牌は字牌役牌より重要}。重なった速さと飜数からいうと、役牌はオタ風より重要。つまり、ほとんどの場合は、「19牌＞役牌＞オタ風」という重要度順でいっていい。例え手牌\ref{mj:yakuhai-0}のようなドラ1の手にとって、立直による打点上昇の幅が大きいので、できるだけ立直を目指す。ブロックが足りていないので、\mj{1m} へのくっつきが \mj{7z} の縦よりうれしい。よって、ここは \mj{1m} ではなく \mj{7z} を切る\citep[][pp.~14--15]{zero-teyaku}。
%
%\tehai{18m 23 667p 557999s 7z}[南家　1巡目　ドラ \mj{3p}][yakuhai-0]

%役牌が頭にするとピンフが消滅するので、ピンフが見える場合は役牌を先に切るべき。たとえ\ref{mj:keepyaku-1}の牌姿、数牌が連続して頭がない。\mj{5z} の縦引きより他の数牌の縦引きが嬉しいのでここで \mj{1s} ではなくて \mj{5z} を切る。
%
%\tehai{27m 2348p 145667s 5z -1p}[東2局　北家　ドラ \mj{8m}][keepyaku-1]

これらのテーマで整理してみると、19牌と字牌の重要度順が下記の３つになる、

\begin{enumerate}
\item 門前が厳しい・リーのみ濃厚な手：役牌＞オタ風＞19牌
\item そこそこの手・立直本線：19牌＞役牌＞オタ風
\item タンピン系手牌・ヘッドレスの手牌：19牌＞オタ風＞役牌
\end{enumerate}


\subsubsection*{役牌より数牌を先に切る場面}

孤立の \mj{12m} と役牌の比較の基準は下記の通り\multicitep{yuusee-yakuand1; yuusee-yakuand2; zero-teyaku, pp.~14--15}：

\begin{itemize}
  \item 役牌を \mj{1m} より先に残す場面：
  \begin{itemize}
    \item あまりにも門前が厳しすぎる（目安：両面と面子が合計で２つ以下）
    \item \mj{14m} のような孤立の筋牌が手牌の中にある
    \item ブロックが足りない、頭がある
    \item ブロックが足りていても愚形がある
    \item ピンフが薄い
    \item \mj{1m} はドラに絡まっていない
  \end{itemize}
  \item 役牌を \mj{2m} より先に残す場面：
  \begin{itemize}
    \item 序盤0面子、かつ、愚形搭子が二つ以上
    \item 端牌がドラ
  \end{itemize}
\end{itemize}

19牌にくっついてできた愚形搭子がうれしくないなら、19牌は孤立牌としての価値が役牌より低くなる。他の部分に愚形搭子があるなら、19牌にくっついてでもどうせ愚形になって、役牌が重なる方より遅い。孤立の19牌がタンヤオのボトルネックだった場合、どうせ払う上、重なった時も役牌より価値が低い。

手牌 \ref{mj:yakuhai-1} にはできた面子がないため、仕掛けた方が遥かに早い。ここは \mj{9m1s} のどちらを切る\citep{yuusee-yakuand1}。

\tehai{1249m 078p 156s 1557z}[東2局　南家　1巡目　ドラ \mj{8s}][yakuhai-1]

手牌 \ref{mj:yakuhai-2} は4ブロック。愚形搭子なし。牌効率から見ると、\mj{9p0s} へのくっつきは \mj{56z} の縦引きより早い。打点から見ると、\mj{9p} へくっついたら \mj{0s} が使いづらくなってしまい、嵌張や辺張の愚形ドラ1になる可能性が高い。頭があるので、役牌が重なった後の仕掛けも見込む。よって、ここは一番不要な \mj{9p} を切る\citep{yuusee-yakuand1}。

\tehai{45677m 239p 0777s 56z}[東1局　東家　3巡目　ドラ \mj{5m}][yakuhai-2]

手牌 \ref{mj:yakuhai-3} の \mj{1p1s} のそばに搭子がある。節~\ref{subsec:lonetile} にも言及されて、\mj{1s} の裏目 \mj{2s} が \mj{45s} に取り上げられない。\mj{1p} がなくなっても裏目の \mj{2p} と \mj{46p} で両嵌張ができるので、ここは裏目が一番痛くない \mj{1p} を切る\citep{yuusee-yakuand1}。

\tehai{124789m 146p 145s 1z}[東1局　南家　4巡目　ドラ \mj{7m}][yakuhai-3]

手牌~\ref{mj:yakuhai-4} の場合はドラが2枚あるけどブロックが4つ、そして \mj{68p} の嵌張があるので、ここは1枚切れの \mj{1z} を残して \mj{1s} を切る。

\tehai{34067m 1168p 16s 167z}[南1局　東家　ドラ \mj{6s}　\mj{1z}が1枚切れ][yakuhai-4]

同じ \ref{mj:yakuhai-5} の牌姿だが、NAGAの選択がドラによる。ドラが \mj{7s} になったらNAGAが \mj{6z} を切る。でなければNAGAは \mj{9s} を切る。

\tehai{44m 45p 1109s 44467z -7m}[東4局　北家][yakuhai-5]

手牌 \ref{mj:yakuhai-6} には面子も両面もない。\mj{89s} はできた辺張だが、このリーのみ濃厚の手牌にとってはうれしくない愚形であろう。\mj{89s} の辺張を払えばタンヤオか役牌かが見えるので、ここは \mj{9s} から切る\citep[][p.~15]{zero-teyaku}。

\tehai{4688m 22p 35589s 167z}[南家　1巡目　ドラ \mj{8p}][yakuhai-6]

手牌 \ref{mj:yakuhai-7} には面子も両面もないので \mj{56z} を切らずに \mj{2m2s} を切る\citep{yuusee-yakuand2}。

\tehai{268m 455899p 279s 56z}[東4局　西家　3巡目　ドラ \mj{1p}][yakuhai-7]

端牌がドラの時、役牌が重なって鳴いてスピードと上がり率を上げる方がいい。よって、\ref{mj:yakuhai-8} は \mj{8s} 切り\citep{yuusee-yakuand2}。

\tehai{2367m 11699p 348s 56z}[東1局　西家　3巡目　ドラ \mj{1p}][yakuhai-8]

一方、同じ三向聴の \ref{mj:yakuhai-9} は\citet[p.~15]{zero-teyaku}によると打 \mj{5z} になる。「手牌がかなり整っている」というのは、おそらく面子が１つできていることを指している。

\tehai{18m 23667p 557999s 5z}[南家　1巡目　ドラ \mj{9s}][yakuhai-9]


\subsubsection*{役牌アタック}

「役牌アタック」というのは、他家が役牌を重ねる前に、もしくは手牌がまた不安定の状態の中で、役牌をオタ風の先に切るという戦術\citep[][pp.~16--17]{zero-teyaku}。役牌が先に出ていくのは、自分の手牌が相当に整っている、もしくは役牌の価値が手組みに対する価値が低い。つまり、

\begin{itemize}
  \item 連続系の形があっても、頭がない。役牌が重なっても鳴いたら手牌が進まない
  \item 字牌以外は全てタンヤオ牌
  \item 門前立直に近い、かつ、ドラ1やドラ2。門前に近いに見える手牌は：
  \begin{itemize}
    \item 両面と面子が合計で３つ以上\citep[][p.~17]{zero-teyaku}
    \item トイツ1つ＋好形2ブロック\citep{senba-yakuhai}
  \end{itemize}
\end{itemize}

上記の条件の一つが満たされば、オタ風より役牌を切ってもいい。たとえ \ref{mj:yakuhai-attack-1} のようなタンヤオ系の手牌と \ref{mj:yakuhai-attack-2} のような門前立直に近い手牌、両方から役牌を先に切ればいい。

\tehai{33468m 22268p 08s 25z}[東場　西家　1巡目][yakuhai-attack-1]

\tehai{1156m 23p 1779s 1347z}[東1局　北家　1巡目　ドラ \mj{5m}][yakuhai-attack-2]


\subsubsection*{1枚切れの字牌と生牌字牌}

\citet[p.~200]{suphx}によると「場1字牌＞19牌＞生牌字牌」。場に1枚切れの字牌の安全度が生牌の方より高いので、序盤で手役に絡まる手牌、アガリ自体の価値が高い役なしの手牌のなど、生牌字牌の攻撃性能が求められる場合以外は、安全度からみて場1の字牌を優先する。生牌字牌を意識的に残そうとする例外は、

\begin{itemize}
  \item 一色手や対子手
  \item 高打点や南場のトップ目で役なし
  \item ブロックが足りていて、数牌も不要であり、引っ張った役牌が別に鳴かれてもいい場合
  \item 守備力がある前提で、
  \begin{itemize}
    \item 役牌の重なりによって絶妙に間に合いそうな手
    \item 役牌対子手
  \end{itemize}
  \item 守備力のないノミ手で気休め程度でも守備力を取る場合
\end{itemize}

立体図~\ref{rittai:jihai-1} は場1の \mj{7z} と生牌の \mj{1z} の二択。ドラも供託もないので、この手牌は仕掛けると1000点に過ぎない。Suphxは千点で仕掛けるために生牌役牌を重ねるようとすることに対しては消極的。それに、上家の捨て牌は変則なので、安全度の高い牌を残す価値が上がっている\citep[][pp.~192--193]{suphx}。

\begin{rittai}[h]
  \centering
  \drawrittai{
    jicha/seat = 北,
    jicha/hand = 45m 11557p 2667s 17z -1p,
    jicha/river = 1m,
    jicha/point = 22100,
    %
    toimen/river = 6z,
    toimen/point = 25000,
    %
    kami/hand = XXXXXXXXXX -66'6z,
    kami/river = 4s 6m,
    kami/point = 24000,
    %
    shimo/river = 7z,
    shimo/point = 28900,
    %
    kyoku = 東1局,
    honba = 1,
    dora = 7m
  }
  \caption{場1の \mj{7z} と生牌の \mj{1z} の二択}
  \label{rittai:jihai-1}
\end{rittai}

立体図~\ref{rittai:jihai-2} では \mj{1z} が重なれば打点が3900から7700まで上昇する。けれどここはトップ目、ブロックも足りているのでSuphxは安全度を優先して \mj{1z} を切った\citep[][p.~194]{suphx}。

\begin{rittai}[h]
  \centering
  
  \drawrittai{
    jicha/seat = 西,
    jicha/hand = 1134488s 134 -4z -555'z,
    jicha/river = 6m 65p,
    jicha/point = 47800,
    %
    toimen/river = 33^s 2^p 9s,
    toimen/point = 7200,
    %
    kami/river = 9^s 3z 8^m 6s,
    kami/point = 13500,
    %
    shimo/river = 3z 2s 7z,
    shimo/point = 31500,
    %
    kyoku = 東4局,
    honba = 2,
    dora = 1m
  }
  
  \caption{2枚見えの \mj{3z} と生牌の \mj{1z} の二択}
  \label{rittai:jihai-2}
\end{rittai}

立体図~\ref{rittai:jihai-3} では重なれば$2000\to7700$もしくは12000になるドラの \mj{4z} だが、一向聴かつすぐ中盤なのでSuphxは \mj{4z} を手放した\citep[][p.~195]{suphx}。これも一向聴の時手牌をスリム化して打点より安全度を優先する打筋。

\begin{rittai}[h]
  \centering
  
  \drawrittai{
    jicha/seat = 北,
    jicha/hand = 34m 2205p 78s 47z -9s -1'11z,
    jicha/river = 1m 1p 2s 8m,
    jicha/point = 25000,
    %
    toimen/river = 1p 17m 3z 7p,
    toimen/point = 25000,
    %
    kami/river = 1m 6z 7m 4s,
    kami/point = 22100,
    %
    shimo/river = 73^z 2m 1z 3s,
    shimo/point = 27900,
    %
    kyoku = 東1局,
    honba = 1,
    dora = 4z
  }
  
  \caption{1枚見えの \mj{7z} と生牌ドラの \mj{4z} の二択}
  \label{rittai:jihai-3}
\end{rittai}

\textbf{打点もないバラバラな手牌は生牌字牌を残さない}\citep[][p~.197]{suphx}。立体図~\ref{rittai:jihai-4} はアガリの可能性がほとんどない手牌。わざと役牌を重ねて上がり率を上げようとすることはない。Suphxは生牌の \mj{7z} を切った。

\begin{rittai}[h]
  \centering
  
  \drawrittai{
    jicha/seat = 北,
    jicha/hand = 23m 3p 1245s 234567z -5p,
    jicha/river = 9m 8p,
    jicha/point = 24400,
    %
    toimen/river = 3z 1m 5z,
    toimen/point = 17400,
    %
    kami/river = 4z 1m 5^z,
    kami/point = 24400,
    %
    shimo/river = 1s 1p 5z,
    shimo/point = 33800,
    %
    kyoku = 東3局,
    dora = 5m
  }
  
  \caption{5向聴かつバラバラの手牌}
  \label{rittai:jihai-4}
\end{rittai}

立体図~\ref{rittai:jihai-5} では仕掛けて7700の決定打。Suphxは2枚切れの \mj{7z} を切らず生牌の \mj{6z} を手放した\citep[][p.~198]{suphx}。いくら一歩も早く上がりたくても、トップ目かつ4向聴の手牌なので、Suphxは安全度を優先した。

\begin{rittai}[h]
  \centering
  
  \drawrittai{
    jicha/seat = 東,
    jicha/hand = 116m 1206p 34s 1567z -1m,
    jicha/river = 9m 23z 8p,
    jicha/point = 34000,
    %
    toimen/river = 9p 9s 1m 7z,
    toimen/point = 22700,
    %
    kami/river = 1s 43z 9^p,
    kami/point = 18300,
    %
    shimo/river = 57z 89^m,
    shimo/point = 25000,
    %
    kyoku = 東4局,
    dora = 5p
  }
  
  \caption{仕掛けて7700の手牌}
  \label{rittai:jihai-5}
\end{rittai}


\subsubsection*{孤立ドラ役牌の扱い方}

\textbf{他の役があれば先に処理する。でなければ他の字牌すら安易に切らない}\citep{yoteru-lonedorayakuhai}。たとえ\ref{mj:lonedorayaku-1}のようなバラバラの配牌、\mj{5z} が重ならなければほとんど上がらない上、切って鳴かせたら損になる。守備力にも構えば字牌は安牌として残して、打牌候補は \mj{8p125s} などになる。

\tehai{24668m 8p 12578s 125z}[ドラ \mj{5z}][lonedorayaku-1]

逆に\ref{mj:lonedorayaku-2}のようにタンヤオ平和赤などな長所があればドラ \mj{5z} が不要になり、他家の \mj{5z} が重なる前に第一打として切るのは正解。

\tehai{34m 468p 334056s 125z}[ドラ \mj{5z}][lonedorayaku-2]

手牌\ref{mj:lonedorayaku-3}はメンタンピン一向聴だが厄介なものを掴まえてしまった場合。ドラを手放す方がいい。ドラがポンされたら捨て牌や場況で対応できるし、ポンされなくて代わりに立直が来る場面でも安牌があるので戦いやすい。

\tehai{23467m 34588p 56s 3z -5z}[7巡目　ドラ \mj{5z}][lonedorayaku-3]

混一色に絡まる場合なら満貫が既にあったらドラを手放す方がいい。手牌\ref{mj:lonedorayaku-4}はもう満貫確定なので、断ラスなど点数が切実に必要な場面でなければ上がり率を優先する。

\tehai{11340899s 125z -7'77z}[東1局　南家　7巡目　ドラ \mj{5z}][lonedorayaku-4]


%---------------------------------------------
\subsubsection{タンヤオ}

\subsubsection*{6ブロックから么九牌のブロックを落とす}

\textbf{6ブロックから么九牌のブロックを落としてタンヤオを確定するのがほぼ正解}\multicitep{gendai-tanyao; zero-teyaku, p.~22}。手牌\ref{mj:tanyao-1}は \mj{46s} のどちらを落とせば受け入れが最大16枚になる。しかし、そのうちピンフがつくツモは \mj{5s} だけ。一方、\mj{9s} を落とせば目前の受け入れが12枚になって4枚損だが、タンヤオが確定されて仕掛けも効くので、速度は実際 \mj{46s} の方より早い。打 \mj{9s}。

\tehai{23467m 345p 446699s}[][tanyao-1]

手牌\ref{mj:tanyao-2}も打 \mj{9m}。牌効率の通りなら打 \mj{2p} だが、他の部分に両面が一つだけ、嵌張待ちになる可能性が高い。打 \mj{9m} とすると、仕掛けが効き、嵌張が頭になる道も解放される。

\tehai{35799m 2467p 24888s}[東1局　南家　3巡目　ドラ \mj{8s}][tanyao-2]

手牌\ref{mj:tanyao-3}も打 \mj{9m}。牌効率の通りなら打 \mj{3m} だが、メンピンで終わる可能性が高い。ドラ引きや345の三色を備えて、打 \mj{9m} の方は打点が高い。

\tehai{356789m 346788p 45s}[東1局　南家　3巡目　ドラ \mj{2m}][tanyao-3]

手牌\ref{mj:tanyao-4}は手なりで打 \mj{7s} なら受け入れが最大の28枚になるが、両面が先に埋められて双碰リーのみになりやすい。ここは一旦 \mj{2p} を切って、2ヘッドの形を維持しつつ、ソーズの部分が伸びたら \mj{1p} を全て払ってもいい\citep[][pp.~68--69]{fkt}。

\tehai{56m 11256p 2246778s}[4巡目　ドラ \mj{6z}][tanyao-4]

手牌\ref{mj:tanyao-4a}はブロックが既に整っている。手なりで立直・ドラ1になりそう。今回の牌姿は四連形と浮いている尖張牌があるため、もう一つのブロックを作るのは容易い。よって、打 \mj{3z} としてタンヤオがほぼ確定されたら向聴数を戻しても構わない\citep[pp.~70--71]{fkt}。

\tehai{3406m 7p 2346688s 33z}[東場　北家　4巡目　ドラ \mj{1p}][tanyao-4a]

タンヤオを狙わない場面は下記の二つ\citep[][p.~23]{zero-teyaku}：

\begin{itemize}
\item 好形搭子が十分
\item タンヤオにすると愚形聴牌になる可能性が高い
\end{itemize}

手牌\ref{mj:tanyao-5}は打 \mj{1m} にすると、両面が先に埋まるとカン \mj{3s} になってしまう。両面両面一向聴が十分なので、ここは牌効率の通りに打 \mj{2s}。

\tehai{1156m 34678p 24666s}[東1局　南家　3巡目　ドラ \mj{8p}][tanyao-5]

手牌\ref{mj:tanyao-6}は牌効率の通りに打 \mj{48p}。両面搭子が十分なので、リャンカンの部分を払うと門前聴牌が無理なく狙える。一盃口を見て \mj{8p} を切るか危険な \mj{4p} を先に切るかが場況次第。

\tehai{4599m 234468p 3478s}[東1局　南家　3巡目　ドラ \mj{8s}][tanyao-6]

手牌\ref{mj:tanyao-7}は牌効率の通りに打 \mj{3m}。手牌\ref{mj:tanyao-2}に似ているドラ3の機会手が、ここは大人しく牌効率を従う。

\tehai{35799m 2367p 23888s}[東1局　南家　3巡目　ドラ \mj{8s}][tanyao-7]

手牌\ref{mj:tanyao-8}もタンヤオを狙わない。\mj{13s} を払えばタンヤオが確定されたが、\mj{25s} の両面が消える上ドラ受けもなくなる。愚形の部分は \mj{68m05p} と \mj{0577p} の受け入れが同じ4枚だが、両面変化の枚数は前者は \mj{5m} 4枚だけ。それに対して後者は \mj{468p} 12枚。よって、ここは嵌張の \mj{68m} を払う\citep[][pp.~32--33]{fkt}。

\tehai{23468m 0577p 12334s}[5巡目　ドラ \mj{5s}][tanyao-8]


\subsubsection*{ダブ東とタンヤオの決断}

\textbf{門前ならダブ東を見切る}\citep[][pp.~138--139]{fkt}。手牌\ref{mj:tanyao-9}はダブ東のために両面搭子を払うのは見合わない。立直・タンヤオ・ドラの方でツモや裏ドラなどで12000になる可能性があるから。

\tehai{22456m 22240p 67s 11z}[東場　東家　6巡目　ドラ \mj{3z}][tanyao-9]

\textbf{仕掛けが入ったらダブ東で12000を目指す}\citep[][pp.~140--141]{fkt}。手牌\ref{mj:tanyao-10}は立直や裏ドラなどがないため、両面両面で行くと2900や都合よくて赤ドラ引いて5800。一方、\mj{1z} を鳴いたら打点が2900から11600になるので段違いです。

\tehai{22456m 40p 67s 11z -6'66z}[東場　東家　6巡目　ドラ \mj{3z}][tanyao-10]

手牌\ref{mj:tanyao-11}のようなダブ東バックでもクイタンよりいい。打点から見ても \mj{1z} の安全度と連打できることから見ても。場況によって打 \mj{3p} でも打 \mj{67s} でもあり\citep[][pp.~143--144, p.~218]{fkt}。

\tehai{22456m 30p 67s 11z -2'22pz}[東場　東家　6巡目　ドラ \mj{3z}][tanyao-11]

%---------------------------------------------
\subsubsection{ピンフ}

\subsubsection*{ドラなしピンフの手筋}

\textbf{ドラなしならリーのみを打たないように両面を固定する}\citep[][pp.~200--201]{doi-prob}。手牌\ref{mj:pinfu-1}は\mj{3m8s} を切ればピンフが確定されている。ここで \mj{3z} をツモった。牌効率の通りならこの \mj{3z} がいらない。しかし、5巡目で自分の手の中にドラがない事実から、他家の手牌の中にドラが通常より多いのが推測できて、先切り・守備力・打点などの要素から考えると、リーのみを避けたいためここで牌効率に従わない。打 \mj{3m}。

\tehai{334789m 56p 11788s -3z}[東1局　西家　5巡目　ドラ \mj{8p}][pinfu-1]

似たような牌姿だが、今回\ref{mj:pinfu-2}にタンヤオが見込むので、タンヤオの優先度が安牌より高い。\mj{788s} を最大限にタンヤオ牌からできるように構う前提で、より危険な \mj{3m} を先切るか、安牌の \mj{3z} を切るかが場況次第。

\tehai{334678m 56p 22788s -3z}[東1局　西家　5巡目　ドラ \mj{8p}][pinfu-2]


%---------------------------------------------
\subsubsection{一盃口}

一盃口は中ぶくれや亜両面などの連続形から自然にできることが多くて、無理に狙う場面が少ない。\mj{33445p} のような両面＋面子の形でも半分の確率で一盃口が消える。\mj{33455p} のような一盃口が確定している形は、自分が待ちの牌をすでに1枚持っているので、もう場に1枚切れになると苦しくなる。両面より \mj{33455p} を先に処理する。2枚見えなら他の嵌張より先に切る\citep{gendai-ipk}。

一盃口を意識していい階段は孤立牌を選ぶ階段とリーのみ濃厚の階段。後者の討論は形次第。前者の討論はドラ枚数次第。手牌\ref{mj:ipk-1}がドラなしなら一盃口で打点上昇のために打 \mj{8p}。ドラ3枚ならダマテンに構えられるために一盃口を狙って打 \mj{8p}。ドラが1枚・2枚なら単純に受け入れの残り枚数で打 \mj{2m} \citep[][p.~32]{zero-teyaku}。

\tehai{12236m 2348p 35799s}[\mj{1m} 1枚切れ][ipk-1]

手牌\ref{mj:ipk-2}はヘッドレス。\mj{1223m} を頭に固定する確率が高くて、一盃口は狙いづらい。打 \mj{8p}。

\tehai{12237m 2348p 23079s}[\mj{1m} 1枚切れ　ドラ \mj{9s}][ipk-2]


\subsubsection*{両面嵌張にするかにいにい形にするか}

\mj{3556778p} という形から \mj{3p} を切るとにいにい形\footnote{\citet{chitist-niinii}がつけた名前。構成が「2枚・1枚・2枚・1枚」から。}になり、\mj{5p} を切ると両面嵌張になる。どちらを切ると受け入れが同じだが、手牌の役・牌の危険度によって最善の一打が変わる。

\begin{itemize}
  \item 両面嵌張
  \begin{itemize}
    \item 面子化する時必ず順子＋順子になる
    \item 危険な内側が先に出る
  \end{itemize}
  \item にいにい形
  \begin{itemize}
    \item 一盃口目がある
    \item 2ヘッドで柔軟
  \end{itemize}
\end{itemize}

\textbf{多くの場合はにいにい形にするのは正解}。手牌\ref{mj:ipk-3}は打 \mj{3p} としたらツモ \mj{5p} や \mj{3s} でテンパネして、ツモ \mj{2s} で \mj{69p} と \mj{14s} を選べることができる。\mj{9p} が急にたくさん切られてくるという緊急事態にも柔軟に対応。赤 \mj{5p} の受けも\citep[][p.~35]{zero-teyaku}。

\tehai{23m 3556778p 33s 777z}[][ipk-3]

そういったメリットは手牌\ref{mj:ipk-4}にはない。役もドラもない手牌、なるべくリーのみを避けたいなら、一盃口が確定されていないにいにい形より、ピンフが確定されている両面嵌張の方がいい。さらに \mj{3p} が \mj{1p} より危険。こう見ると、危険な \mj{3p} を先に切る\citep[][p.~35]{zero-teyaku}。

\tehai{23m 1334556p 33789s}[ドラ \mj{7m}][ipk-4]

手牌\ref{mj:ipk-5}も似ている形の牌姿。今回の牌姿は既にドラ1があって、「なんとなく役をつくる」という念頭がないし、\mj{3m5m} も危険度に差がない。そうすると、打点最大値を見て打 \mj{3m} として一盃口を狙う\citep[][pp.~120--121]{uzaku-haikouritsu}。ただし、トップ目でピンフでダマしたい場面などなら打 \mj{5m} もあり。

\tehai{3556778m 45p 12388s}[東家　11巡目　ドラ \mj{2s}][ipk-5]

\begin{table}[ht]
  \centering
  \begin{tabular}{lll}
    切る牌 & \mj{3m} & \mj{5m} \\
    \hline
    受け入れ & \mj{569m36p8s} 19枚 & \mj{469m36p} 19枚 \\
    ピンフツモ & \mj{69m36p} (78.9\%) & \mj{469m36p} (100\%) \\
    一盃口ツモ & \mj{6m} (15.8\%) & (0\%)
  \end{tabular}
  \caption{手牌\ref{mj:ipk-5}の受け入れ}
  \label{tab:ipk-5}
\end{table}


\subsubsection*{三連飛び対子}

\mj{113355p} のように連なっている形は「三連飛び対子」\citep[][p.~35]{zero-teyaku}や「2回飛んだ飛び対子」\citep{yoteru-tobitoitsu}と呼ばれている。多くの場合は、三連飛び対子をもう一つの搭子とまとめてどちらを払うのを考える\citep[][pp. 114--117]{uzaku-haikouritsu}。表~\ref{tab:santobitt-1} の通り、搭子が対子でない限り、三連飛び対子から真ん中を抜く方は受け入れが多い。けれど、そうするとツモ \mj{24m} でも一盃口ができない。他の搭子は愚形で、リーのみになりやすいし待ちとしても強くないならそのデメリットはひどい。

\begin{table}
  \centering
  \begin{tabular}{lll}
    形 & 打 \mj{3m} の受け入れ & 打 \mj{8s} の受け入れ \\
    \hline
    \mj{113355m78s} & \mj{1245m69s} 20枚 & \mj{12345m} 14枚 \\
    \mj{113355m68s} & \mj{1245m7s} 16枚 & \mj{12345m} 14枚 \\
    \mj{113355m89s} & \mj{1245m7s} 16枚 & \mj{12345m} 14枚 \\
    \mj{113355m88s} & \mj{1245m8s} 14枚 & \mj{12345m} 14枚 \\
  \end{tabular}
  \caption{三連飛び対子＋搭子から1枚を外した後の受け入れ}
  \label{tab:santobitt-1}
\end{table}

例え手牌\ref{mj:santobitt-1}は \mj{113355m99p} から \mj{3m} と \mj{9p} の二択。それぞれのメリットは、

\begin{itemize}
  \item 打 \mj{3m}
  \begin{itemize}
    \item \mj{1m9p} の双碰になるとアガリ率が高い
  \end{itemize}
  \item 打 \mj{9p}
  \begin{itemize}
    \item 一盃口目が残る
    \item 連打できるので安牌やフォロー牌が持てる
  \end{itemize}
\end{itemize}

「双碰として強い」というメリットは、\mj{9p} が場に1枚見えになると消える。そうすると \mj{9p} を切って一盃口を狙う\citep[][p.~35]{zero-teyaku}。

\tehai{113355m 123 99p 567s}[][santobitt-1]

手牌\ref{mj:santobitt-2}の三連飛び対子の部分が中寄りになって、\mj{2m} 双碰が待ちとして弱くなる。ここは \mj{1p} 切りに寄る\citep{yoteru-tobitoitsu}。

\tehai{224466m 11p 789s 123s}[][santobitt-2]


\subsubsection*{一盃口クラッシュ}

一盃口は門前役で、鳴いたら一盃口が消える。が、一盃口の形は頭が作りやすくて、仕掛け手順も意外にも存在する。手牌\ref{mj:ipk-crash}から単騎を取らずに \mj{4p} を落とすという手順は「一盃口クラッシュ」と呼ばれている。上がりづらい単騎より、ポン \mj{35p} とチー \mj{2356p} で聴牌できる広い一向聴をとるべき\citep[][p.~36]{zero-teyaku}。

\tehai{67m 334450p 06s -77'7z}[ドラ \mj{5m}][ipk-crash]


%---------------------------------------------
\subsubsection{混一色・清一色}

混一色を狙う場面は下記の通り、

\begin{itemize}
\item 受け入れが少なくて、愚形が多くて、門前聴牌まで遠い\citep{shibukawa-honitsu}
\item 鳴いても満貫が見える\citep{hirasawa-honitsu}
\item 手なりで立直打っても2飜以下\citep{hirasawa-honitsu}
\item 役牌対子が二つある（すなわち役役）\citep[][pp.~96--97]{zero-teyaku}
\item 役牌対子が一つ、他の部分がバラバラ\citep[][pp.~97--98]{zero-teyaku}
\end{itemize}

混一色は、たった一つの色の数牌と字牌しか受け入れないので、牌効率としては非効率的。しかし、仕掛けが効く上、字牌を持っていくことで守備力も十分で、打点も十分。特に立直から遠い悪い手牌こそ、混一色で守備力と打点を維持して進むのがオススメ。


\subsubsection*{役役の場合は混一色を狙う}

麻雀は打点効率として満貫を作るのは効率が一番高い。役牌対子が二つある場合、両飜の役がつけば満貫になるので、混一色やトイトイの価値が高まる。そして混一色の方が受け入れが多いので、役役の場合は常に混一色を意識する\citep[][pp.~96--97]{zero-teyaku}。よって、手牌\ref{mj:honitsu-1}は染めていく。打 \mj{8s}。

\tehai{7m 12456p 228s 44557z}[東1局　北家　1巡目　ドラ \mj{3z}][honitsu-1]

手牌\ref{mj:honitsu-2}はソーズの部分が完成している。が、超序盤（4巡目以内）で役牌が出たらポンして混一色を狙ってソーズの面子を落とす\citep[][p.~97]{zero-teyaku}。

\tehai{12456p 123s 44557z}[東1局　北家　配牌　ドラ \mj{3z}][honitsu-2]


\subsubsection*{役牌対子がない場合は基本的に狙わない}

手牌\ref{mj:honitsu-3}はすでに頭と一盃口が整っていて、\mj{8p26s} にくっつけば立直・タンヤオ・一盃口で打点が十分。一方、混一色を狙って仕掛けたら概ね3900。僅かな清一色の可能性を残して、\mj{37z} から切る\citep{shibukawa-honitsu}。

\tehai{144556688m 8p 26s 37z}[南1局　南家　1巡目　ドラ \mj{5p}　ダンラス][honitsu-3]

手牌\ref{mj:honitsu-4}には5枚の孤立字牌があるが、字牌を重なるのに待つより、数牌のくっつきでタンヤオやピンフを狙う方が早い。ここは字牌を処理する\citep{yuusee-honitsu}。

\tehai{48m 23457p 25s 12456z}[東1局　西家　1巡目　ドラ \mj{1m}][honitsu-4]

手牌\ref{mj:honitsu-5}と\ref{mj:honitsu-6}はともに受け入れが少ない愚形だらけの手牌、門前聴牌は遠く見えるので、守備力のために消極的に染めていく\citep{shibukawa-honitsu}。

\tehai{114478m 246s 44567z}[東3局　東家　2巡目　ドラ \mj{2m}　\mj{4m} 1枚切り][honitsu-5]

\tehai{35m 57p 123449s 2244z}[東3局　東家　2巡目　ドラ \mj{8s}　\mj{4z} 1枚切り][honitsu-6]

手牌\ref{mj:honitsu-yakuyaku-1}にも6ブロックがある。対子が四つなので満貫を作るためトイトイも視野に入る。だとすると一番落としたい部分は \mj{56m} の両面と \mj{35z} の孤立牌。\mj{3z} は自風なので \mj{5z} が先。赤ドラがある場合なら赤受けで \mj{6m} が先。手なりの2000で上がる手順を残したいなら \mj{6z}。混一色とトイトイの可能性を最大限に高めるなら \mj{6m} \citep[][pp.~98--103]{tetsusenryaku}。

\tehai{56m 22p 2499s 113566z}[東1局　西家　1巡目　ドラ \mj{8p}][honitsu-yakuyaku-1]

手牌\ref{mj:honitsu-yakuyaku-2}は\ref{mj:honitsu-yakuyaku-1}に似ているが、マンズの部分が嵌張になり、ソーズの下が両嵌張になる。こうなるとマンズの部分が明確に弱くになって、\mj{46m} を切って混一色に行く\citep[][pp.~98--103]{tetsusenryaku}。

\tehai{46m 22p 24699s 11566z}[東1局　西家　1巡目　ドラ \mj{8p}][honitsu-yakuyaku-2]


\subsubsection*{役牌対子が一つの場合は他の部分次第}

役牌対子が一つ、孤立の役牌が二つなら、役牌の重なりの価値が高まる。これは「役役ダッシュ準備状態」と言い、他の色の孤立の1289を切る\citep[][p.~97]{zero-teyaku}。よって、手牌\ref{mj:honitsu-7}は打 \mj{8p}、\ref{mj:honitsu-8}は打 \mj{8s}。

\tehai{1126m 8p 56889s 1566z}[東1局　北家　1巡目　ドラ \mj{2s}][honitsu-7]

\tehai{34m 1455689p 8s 4556z}[東1局　北家　1巡目　ドラ \mj{2m}][honitsu-8]

手牌\ref{mj:honitsu-9}も打 \mj{28s} でピンズの混一色を狙う\citep{hirasawa-honitsu}。

\tehai{347m 34479p 28s 1157z}[ドラ \mj{1p}][honitsu-9]

手牌\ref{mj:honitsu-10}は混一色を狙わない。孤立の役牌が一つ、そしてマンズの部分は伸びにくい。ここは素直に \mj{6z} を切る\citep[][p.~98]{zero-teyaku}。

\tehai{35999m 379p 48s 5556z}[東1局　北家　1巡目][honitsu-10]

手牌\ref{mj:honitsu-11}も混一色を狙わない。他の色の部分が全て両面なので、立直が近くなる。

\tehai{67m 23p 223056s 5567z}[東1局　北家　1巡目 ドラ \mj{1s}][honitsu-11]

手牌\ref{mj:honitsu-12}は他の色に搭子がない。ピンズの部分に一盃口の形があるが、立直しても2600が現実的。混一色に行くなら最低3900、自風の \mj{4z} か重なったら満貫が見える。オタ風の \mj{3z} もワンチャントイトイの種にもなる。赤受けを考えると打 \mj{7s} \citep[][pp.~104--109]{tetsusenryaku}。

\tehai{5m 22333447p 7s 3477z}[東1局　北家　1巡目 ドラ \mj{6z}][honitsu-12]


\subsubsection*{2番目に多い色から切る}

手牌\ref{mj:honitsu-13}は親からの第一打の \mj{4z} をポンしかける場面。\mj{6m46p} のどちらも不要だが、\mj{6m} を遅らせて他家にソーズ寄せを悟られないようにするのは混一色のコツ。この後マンズを引いたらツモ切り\citep[][pp.~118--119]{suzuki}。

\tehai{6m 46p 14467s 33446z -4z}[東場　北家　1巡目　ドラ \mj{2m}][honitsu-13]

手牌\ref{mj:honitsu-14}はダブ東をポンしたところ。ピンズとソーズがどちらも2枚ずつ。枚数が同じだが、外側切ると河が手なりに見えるのでここは \mj{8p} を切る\citep[][pp.~120--121]{suzuki}。

\tehai{1179m 68p 36s 447z -111'z}[東場　東家　2巡目][honitsu-14]


%---------------------------------------------
\subsubsection{役役の進行}

\textbf{役牌の対子が二つあるなら満貫を目指せ}\citep{yoteru-yakuyaku}。「役役」は手牌の中に二つの役牌の対子のことを指し、最低の得点の2000を7700や満貫にすれば四倍になるので打点上昇のスィートスポット。役役にさらに2飜をつく方法は(1)混一色(2)トイトイ。

例え\ref{mj:yy-1}の手牌は嵌張だらけで、\mj{34s} が唯一の両面。\mj{34s} の主線で進行すればただの2000点。うまくいけなくて役牌対子は一つしか刻子にならない場合は1000点。どうせアガりづらくて安いの手牌ならアガリ率を下げて手牌を染めにいくの方が得。よって \mj{68p34s} の全てが打牌候補。

\tehai{2579m 68p 34s 135577z}[東1局　南家　1巡目　ドラなし][yy-1]

手牌\ref{mj:yy-2}ならトイトイを目指す。トイトイなら中張牌が重なっても嬉しくないのでここで \mj{3m} を切る。

\tehai{223m 119p 28s 135577z}[東1局　南家　1巡目　ドラなし][yy-2]


\subsubsection*{ドラが1枚持っている場合}

\tehai{67m 1330p 46s 135577z}[東1局　南家　1巡目　ドラ \mj{2m}][yy-3]

\textbf{ドラが1枚持っているなら手なりで打つのが基本}。無理矢理染めても3900 $\to$ 8000は二倍に過ぎないので効率的に良くない。よって、手牌\ref{mj:yy-3}は染めないで、素直に \mj{13z} から切る。

\tehai{137m 3799p 124s 5577z}[東1局　南家　1巡目　ドラ \mj{1m}][yy-4]

そういっても満貫が見込めるケースは\ref{mj:yy-4}のようにチャンタがついている手牌。他に三色で満貫にするケースも。

\tehai{13m 67s 5577z -6'0799'9m}[東1局　南家　7巡目　ドラ \mj{8p}][yy-5]

鳴いたけど役牌が一つしか刻子にならない牌姿なら同じ、できれば満貫にする。手牌\ref{mj:yy-5}は混一色が見えるので \mj{67s} の両面は割愛。

\tehai{405s 25577z -1'11p888'm}[東1局　南家　7巡目　ドラ \mj{8p}][yy-6]

手牌\ref{mj:yy-6}の場合も \mj{4s} を切ってトイトイで満貫を目指す。もちろん場況・点棒状況によって素直に \mj{2z} を切るケースもあるけど、平場の東1局なら満貫を取る。


\subsubsection*{役役で真っ直ぐに行くケース}

\tehai{2467m 223p 56s 25577z}[東1局　南家　1巡目　ドラ \mj{1m}][yy-7]

\textbf{両面が三つあるなら素直に行く}。手牌\ref{mj:yy-7}のように両面がいっぱいの手牌を無理矢理染めると効率が悪過ぎて打点がいくら高くてもアガらないと意味がない。けれどやはり良形と愚形がともにある手牌は染めるか否かのは紙一重。


%---------------------------------------------
\subsubsection{三色同順}

「配牌を見たら三色を探せ」は、阿佐田哲也氏が残した言葉。昭和時代の麻雀は一発・赤ドラや裏ドラなどの役がまだ普及されていないため、三色などの手役で打点を作るのが大事。平成になると、麻雀ルールがインフレしてきたため、「最速聴牌」が優先され、手役の重要度が落ちて、「三色は偶然役」という意見も。

そう言っても、ここまでの話もしばしば言及したように、全ての場面で「棒テン即リー」は、科学的に局収支から見ると必ずしも正しくない。よりいい結果に辿り着くため、牌効率を逆らって、もしくはもっと聴牌から遠い階段で工夫を入れて手役の可能性を最大限にする必要がある。そうすると三色を狙う技術と基準が重要。三色を狙う基準は\citet{gendai-ssk}によると：

\begin{itemize}
\item 確定しやすさ
\item 確定しづらい場合、三色を狙うことによる受け入れのロス
\item 三色の２翻がつくことによる打点面のメリット
\end{itemize}

手牌にドラはないなら打点を作るため三色が重要になる。そして手牌の価値が非常高いが門前聴牌まで遠い場合、三色による仕掛けのメリットが高い。

手牌が一定の程度で整っているなら三色は狙いづらい。例え手牌\ref{mj:ssk-1}に123の三色部分が7枚もあるが、頭がないので実は123の順子が三つ揃いづらい。この手牌はピンフ狙いの効率が高い。打 \mj{1m}。手牌\ref{mj:ssk-2}も三色部分以外で面子が二つ完成しているので三色を狙わない。打 \mj{1p} \citep[][p.~67]{zero-teyaku}。

\tehai{122m 12357p 236778s}[南家　5巡目　ドラ \mj{2p}][ssk-1]

\tehai{235678m 13567p 123s}[南家　5巡目　ドラ \mj{8m}][ssk-2]


\subsubsection*{リーのみの場合は受け入れより手役}

手牌\ref{mj:ssk-3}と\ref{mj:ssk-4}は打 \mj{7p} で完全一向聴になる。手牌\ref{mj:ssk-3}はそうしべきだが、手牌\ref{mj:ssk-4}はタンヤオが付かないので完全一向聴にするとリーのみになってしまう可能性がある。打 \mj{6m} によって \mj{16m} の受けが消えるが、ツモ \mj{7m6p} で三色が見える。期待値として打 \mj{6m} の方がいい\citep{clearrain-ssk}。

\tehai{22566m 457p 222567s}[6巡目　ドラなし][ssk-3]

\tehai{11566m 457p 222567s}[6巡目　ドラなし][ssk-4]

手牌\ref{mj:ssk-5}は複合搭子二つ、両面対子が一つの形。打 \mj{4m} と打 \mj{6s} の受け入れが同じ。表~\ref{tab:fkgtt-2} の通りなら、両面聴牌の枚数のために打 \mj{6s}。が、この手牌はリーのみの上、打 \mj{4m} と234の三色が確定するので、ここは打 \mj{4m}。両面リーのみより役が大事\citep[][p.~65]{zero-teyaku}。

\tehai{244m 234p 234667s 44z}[南家　3巡目　ドラ \mj{3z}][ssk-5]

手牌\ref{mj:ssk-6}はと同じだが、ドラが2枚なので素直に良形聴牌率で打 \mj{6s} \citep[][p.~65]{zero-teyaku}。

\tehai{244m 234p 234667s 44z}[南家　3巡目　ドラ \mj{4z}][ssk-6]


\subsubsection*{悪い配牌こそ手役を狙うべき}

手牌\ref{mj:ssk-7}は4向聴でバラバラすぎる悪配牌。上がりに速度が必要なので、仕掛けが効く役といえばタンヤオ、役牌、そして三色。今回 \mj{1s} があって123の三色が見込む。打 \mj{9s} \citep[][p.~64]{zero-teyaku}。

\tehai{1356m 124p 15899s 56z}[南家　1巡目][ssk-7]


\subsubsection*{ブロックより三色の種}

手牌\ref{mj:ssk-8}に6ブロックがある。\mj{8s} は孤立牌なので、牌効率通りなら打 \mj{8s}。が、\mj{2244p33z} が3ブロックとして効率が悪いので、ここは打 \mj{2p} として \mj{244p} の1ブロックにして三色の種を残す\citep[][p.~64]{zero-teyaku}。

\tehai{79m 2244789p 458s 33z}[南家　3巡目　ドラ \mj{9p}][ssk-8]

手牌\ref{mj:ssk-9}も \mj{2446m} が2ブロックとして面子を作るのが効率的に悪い。打 \mj{2m} によって1ブロックにして \mj{7p} を残して789の三色を狙う\citep{gendai-ssk}。

\tehai{2446789m 7p 113578s}[][ssk-9]


\subsubsection*{6ブロックの選択}

手牌\ref{mj:ssk-10}は789の三色部分に7枚があるが両面両面なので実は確定していない。そして \mj{8m} を落とすと \mj{233457m} という両面嵌張ができて、牌効率通りの打法。よって、どっちを切るのはドラ枚数次第。ドラがないなら打 \mj{3m} として三色と一通を見て、ドラが1枚以上なら牌効率の通りに打 \mj{8m} \citep{gendai-ssk}。

\tehai{2334578m 789p 5578s}[][ssk-10]

手牌\ref{mj:ssk-11}は6ブロックがある。6ブロック理論によると \mj{4688p} は2ブロックにしないで、\mj{688p} の1ブロックにする方が効率的だが、\mj{34m} 両面固定によって上がり率と三色の可能性が高まる。打 \mj{3m} \citep[][p.~65]{zero-teyaku}。

\tehai{334m 2344688p 3488s}[][ssk-11]

手牌\ref{mj:ssk-12}は三色不確定の形。両面と対子しかないなので、ここは三色に拘らない。受けが被っている \mj{67p} を落とす\citep{gendai-ssk}。

\tehai{06m 34467p 1112367s}[][ssk-12]

%手牌\ref{mj:ssk-13}には567の三色が確定しているが、タンヤオがあるので（平場なら）三色に拘る理由はない。牌効率の通りに打 \mj{5m} \citep{gendai-ssk}。
%ssk-11もタンヤオがあるのになぜssk-11は三色狙う？
%\tehai{5788m 223457p 3357s}[][ssk-13]


\subsubsection*{孤立牌同士の比較}

手牌\ref{mj:ssk-14}は一向聴、すでに \mj{789m789p} があって、\mj{9s} にくっつけば三色が確定する。が、他の孤立牌は両面できやすい \mj{3m3p}。よって、この手牌の打牌は巡目とドラ枚数次第。ドラがないなら \mj{3m3p} を切って789の三色を狙う。ドラが2枚以上あるなら牌効率の通りに打 \mj{9s}。ドラが1枚ならどちらもいい。巡目が早ければ3色を狙い、遅ければ牌効率に従う\citep{yoteru-ssk}。

\tehai{3789m 3789p 229s 333z}[南家][ssk-14]

手牌\ref{mj:ssk-15}は\ref{mj:ssk-14}より聴牌まで遠い手牌。同じ \mj{9s} にくっつけば三色が確定する。ここでドラがないと三色を狙うという方針が\ref{mj:ssk-14}と同じだが、ドラが3枚以上なら、仕掛けを効かせるために三色を狙う価値が高まるので打 \mj{3m} \citep{yoteru-ssk}。

\tehai{3789m 3379p 3559s 33z}[南家][ssk-15]

手牌\ref{mj:ssk-16}にもう4ブロックがある。牌効率として \mj{7s} から両面ができるので \mj{1s} より強いが、\mj{3p} を先に引かなければ嵌張リーのみになってしまう。よって、ここは \mj{1s} を残して打 \mj{7s} \citep{gendai-ssk}。

\tehai{123567m 1255p 17s56z}[][ssk-16]


\subsubsection*{三色の種にも強さの差がある}

\textbf{真ん中の種は弱い}\citep[][pp.~65--66]{zero-teyaku}。手牌\ref{mj:ssk-17}の \mj{3p} は下の三色の種、\mj{7p} はドラなのでどっちも切らない。そうすると \mj{14m} のどっちを切る問題になる。打 \mj{1m} と打 \mj{4m} の比較は表~\ref{tab:ssk-1} の通り、\mj{3p} に \mj{124p} がくっついたら差がでて、特にツモ \mj{1p} の先に打 \mj{1m} の場合は役なしとなり、打 \mj{4m} の場合は三色が確定するなのでカバーできる。打 \mj{4m} で123の三色を狙う。

\tehai{123405m 37p 122334s}[南家　5巡目　ドラ \mj{7p}][ssk-17]

\begin{table}[h]
  \centering
  \begin{tabular}{lll}
    役 & 打 \mj{1m} & 打 \mj{4m} \\
    \hline
    ツモ \mj{1p} & なし & 三色 \\
    ツモ \mj{2p} & (三色)・ピンフ & (三色)・ピンフ \\
    ツモ \mj{4p} & (三色)・ピンフ & ピンフ
  \end{tabular}
  \caption{手牌から \mj{14m} のどちらを切る選択肢の比較。括弧は三色が確定しない場合。}
  \label{tab:ssk-1}
\end{table}

手牌\ref{mj:ssk-18}と\ref{mj:ssk-19}も三色の種がある牌姿。この種を切ると愚形立直ドラ1になる。先のように比較をすると、\mj{7m} にピンフが付かないと三色もない。一方、\mj{6m} にピンフがつくなくても \mj{8m} 引いたら三色が確定する。よって、手牌は打 \mj{7m} として立直。手牌は \mj{13p} を落とす\citep[][p.~66]{zero-teyaku}。

\tehai{1237m 13678p 67778s}[南家　3巡目　ドラ \mj{1m}][ssk-18]

\tehai{1236m 13678p 67778s}[南家　3巡目　ドラ \mj{1m}][ssk-19]


\subsubsection*{両天秤}

「両天秤」は複合しない二つの役の準備状態を指す。どっちの役になるのかがツモ次第。三色と一通の両天秤があり、異なる三色の両天秤もある。手牌\ref{mj:ssk-20}は打 \mj{4m} にすると、\mj{5678m67p67s} の形で \mj{5p} 引いたら打 \mj{8m} で567の三色が見えて、\mj{8s} 引いたら打 \mj{5m} で678の三色が見える。つまり、567と678の三色の両天秤。牌効率の通りなら打 \mj{67p} として受け入れが最大22枚だが、打 \mj{4m} にしても受け入れも19枚あるのでその3枚損のリターンが大きい\citep{rb1}\footnote{もちろん、点棒状況・巡目・ドラ枚数によって三色狙わない場面もある。}。

\tehai{45678m 3367p 34567s}[][ssk-20]

手牌\ref{mj:ssk-21}も123と234の三色の両天秤。打 \mj{3p} は打 \mj{14s} に対して4枚損だが、立直・ドラ・三色の方が立直・ドラ・ピンフより高いので相当なリターン\citep{clearrain-ssk}。

\tehai{23067m 23399p 1234s}[6巡目][ssk-21]

手牌\ref{mj:ssk-22}も345と456の三色の両天秤。今回はリーのみになる可能性があるので、三色の両天秤を狙う価値がもっと高まる\citep{gendai-ssk}。打 \mj{4p}。

\tehai{3456m 445 999p 45 88s}[][ssk-22]

手牌\ref{mj:ssk-23}は打 \mj{2m3p} にすると567と678の三色の両天秤になるが、打 \mj{58s} に対して実は11枚損。暗刻ありヘッドレスの一向聴なので、良形メンタンが確定している。よって、この11枚の損失は大きすぎて、三色の両天秤にしない。567と678の三色のどちらを選ぶと、8の方がでやすくて、678の方が完成しやすいので打 \mj{5s} \citep{clearrain-ssk}。

\tehai{22267m 33367p 5678s}[6巡目][ssk-23]

手牌\ref{mj:ssk-24}にはもうメンタンピンドラ1（＝満貫）なので三色の両天秤にしない。牌効率の通りに打 \mj{8m} \citep{gendai-ssk}。

\tehai{566678m 306778p 67s}[][ssk-24]


%---------------------------------------------
\subsubsection{一気通貫}

一通は三色と同じ、門前で2飜、喰い下がりで1飜。が、一通の完成形に一色の1から9までの数牌が必要なので、組み合わせの数（=3）は三色（=7）より少ない。従って、一通の出現確率は三色より低い。


\subsubsection*{一通が狙える時}

聴牌まで遠い階段、またはドラなしの手牌は、一色の牌は5種類あるなら一通を視野に入れる。孤立牌や愚形搭子を比較する時、一通のために牌効率を逆らうことが多い。手牌\ref{mj:ittsu-1}は孤立の \mj{7m9p8s} の比較。牌効率の通りなら \mj{9p} だが、この手牌はピンフもタンヤオもドラもない。ドラが \mj{6m} なのでドラを使うため \mj{7m} は切らない。そして、\mj{2345679p} に \mj{18p} 引いたら一通が見えるので \mj{8s} を切る\citep[][p.~68]{zero-teyaku}。

\tehai{3457m 2345679p 118s}[ドラ \mj{6m}][ittsu-1]

手牌\ref{mj:ittsu-2}は \mj{1m1s} の比較。マンズの部分は5種類があって、一通にはちょっと遠いが、それでもより孤立の \mj{1s} は先に切るべき。手牌\ref{mj:ittsu-3}も同じ理由で打 \mj{2p} \citep[][p.~69]{zero-teyaku}。

\tehai{14578m 1136p 1899s 5z}[ドラ \mj{5m}][ittsu-2]

\tehai{267899m 2678p 1399s}[][ittsu-3]

手牌\ref{mj:ittsu-4}は辺張 \mj{12m} と嵌張 \mj{13s} の比較。牌効率通りなら、ツモ \mj{4s} の両面変化を見て辺張を落とす。が、マンズの部分が6種類があって、ツモ \mj{45m} による一通チャンスを逃さないようにここは \mj{13s} を落とす\citep[][pp.~68--69]{zero-teyaku}。

\tehai{1267899m 678p 1399s}[][ittsu-4]

手牌\ref{mj:ittsu-5}は辺張 \mj{12m} と辺張 \mj{89m} と両面 \mj{45s} の比較。\textbf{一通確定愚形愚形は、何の役も確定していない愚形＋両面より強い}\citep{gendai-ittsu}。よって、打 \mj{45s}。

\tehai{1245689m 99p 45567s}[][ittsu-5]


\subsubsection*{一通を見切る時}

一色の部分に7種類があっても、1枚・2枚を払えばタンヤオがつく場合、もしくはドラはすでに1枚ある場合、一通を狙うことは滅多にない。例えば、手牌\ref{mj:ittsu-6}にマンズの部分が7種類。が、\mj{9m} を切っても \mj{27m} の受け入れが減らない上、さらに \mj{1m} を切るとタンヤオがつく。正着は打 \mj{9m}。手牌の形を限定する一通は、タンヤオ・ドラ1より弱い\citep[][p.~69]{zero-teyaku}。

\tehai{1345689m 33468p 06s}[][ittsu-6]

手牌\ref{mj:ittsu-7}のようなドラなしタンヤオにもならない手牌なら \mj{26p} を切って一通を見る。ドラが1枚あるならやはり打 \mj{9m} \citep[][p.~70]{zero-teyaku}。

\tehai{1345689m 11246p 78s}[][ittsu-7]

手牌\ref{mj:ittsu-8}から \mj{19m} を外すとタンヤオがあるので一通を見ない。打 \mj{19m} \citep[][p.~70]{zero-teyaku}。

\tehai{1345679m 3557p 455s}[][ittsu-8]

手牌\ref{mj:ittsu-9}から辺張 \mj{12m} を落とすとピンフがつくので一通は見ない。打 \mj{1m} \citep{gendai-ittsu}。

\tehai{1256789m 99p 45567s}[][ittsu-9]


\subsubsection*{黄金の一向聴}

くっつき一向聴の中で、三色と一通の両天秤の牌姿は「黄金の一向聴」と呼ばれている。例え手牌\ref{mj:golden-1}は打 \mj{9s} にすると、ツモ \mj{16p} で一通が見えて、ツモ \mj{79s} で789の三色が見える。打 \mj{8m} のピンフより価値が高いので、ここは打 \mj{9s} \citep{rb1}。

\tehai{2345789p 556789s 8m}[][golden-1]

手牌\ref{mj:golden-2}も黄金の一向聴。黄金一向聴は必ず\ref{mj:golden-1}のように明確に浮いている \mj{8m} があるわけではない。くっつき一向聴なら\ref{mj:golden-2}のような四連形＋両面対子の形にもなりうる。

\tehai{2346789p 234m 344s}[][golden-2]

手牌\ref{mj:golden-3}も打 \mj{6s} による黄金の一向聴。打 \mj{3m} による四連形＋四連形のくっつきの受け入れが53枚で一番広いが、そのうち愚形リーのみになってしまうツモも多い\citep{shibukawa-ssk}。

\tehai{344m 2345789p 3456s}[東1局　東家　5巡目　ドラ \mj{2z}][golden-3]

手牌\ref{mj:golden-4}は打 \mj{6m} とすると345の三色とソーズの一通の両天秤。が、四連形にピンフツモは \mj{2457m} 4種類があって、\mj{3p} へのくっつきの \mj{24p} より多い。この手牌は三色や一通を特に見なくてもメンタンピンも十分なので打 \mj{3p} \citep{gendai-ittsu}。

\tehai{345688m 3p 2345678s}[][golden-4]

ですが、対子の部分を \mj{9m} に変えるとタンヤオが消えるので三色や一通の重要さが大きくなって、打 \mj{6m} の方が優勢になる\footnote{シミュレーションによる結論}。

\tehai{345699m 3p 2345678s}[][golden-5]


\subsubsection*{6ブロックから一通と三色の決断}

三色と一通がともに3ブロックあって、両天秤にならないためどちらを選ぶときは\textbf{三色にすると有利になることが多い}\citep{gendai-ittsu}。一通は確定しづらい点と、タンヤオと複合しやすい点からです。

手牌\ref{mj:ittsu-10}は789の三色とピンズの一通がともに見える。が、この手牌は6ブロック、\mj{78m} と \mj{23p} と \mj{56p} のどちらを落とさなければならない。この手から一通を成立させるには \mj{14p} 両方が必要なので、まず一通は確定しづらい。\mj{78m} を落として \mj{1p} を引いても待ちが \mj{47p} になって、実は弱い。よって、この手牌は一通を狙わない。打 \mj{5p} とすると、ツモ \mj{6s} で打 \mj{9s9p} として678の三色へ移行することにより、タンヤオも見える。

\tehai{3378m 2356789p 789s}[東場　北家　7巡目　ドラ \mj{8s}][ittsu-10]

手牌\ref{mj:ittsu-11}はピンズの一通が確定している。この手から両面を落とす後、他の両面が先に埋まってしまうとカン \mj{8p} の5200になる。一方、打 \mj{9p} により123と234の三色の両天秤にすることができるし、ツモ \mj{4m4p} ならタンピン三色まで狙える。この手牌も一通を狙わないで、打 \mj{9p}。

\tehai{23m 12345679p 2377s}[東場　西家　8巡目　ドラ \mj{3p}][ittsu-11]

%---------------------------------------------
\subsubsection{トイトイ}

トイトイに喰い下がりはないため、ドラや他の役に複合すると5200や8000になりやすい。高打点を作るために相性がいい役だが、立直を受けている時自分の手牌の種類は限定されていて守備力は低下\citep{yoteru-toitoi}。門前や七対子の準備状態からトイトイの移行は「鳴き判断」一章を参照。

\subsubsection*{1300から5200へ}

手牌\ref{mj:toitoi-1}は打 \mj{6m} とすると打点が1300しかない。トイトイにして、聴牌枚数が半分になるが打点が四倍になる。よって、ここは打 \mj{5m} \citep[][p.~57]{zero-teyaku}。

\tehai{566m 22288p 222s -666'z}[ドラ \mj{1z}][toitoi-1]

2副露でもトイトイに受けた方がいい。手牌\ref{mj:toitoi-2}も打 \mj{5m} \citep[][p.~57]{zero-teyaku}。

\tehai{566m 88p 222s -22'2p666'z}[ドラ \mj{1z}][toitoi-2]

両面聴牌に満貫があるならトイトイで跳満にする必要はない。聴牌枚数を半分にしても打点が$8000\to12000$、つまり1.5倍しかないので、ここは打 \mj{6m} \citep[][p.~57]{zero-teyaku}。

\tehai{566m 88p 222s -22'2p666'z}[ドラ \mj{2p}][toitoi-3]

手牌\ref{mj:toitoi-4}はツモ \mj{1s} で4つの対子ができて、トイトイにすると5200になる。ここは \mj{3s7p7m} の順番に切っていく\citep[][p.~57]{zero-teyaku}。

\tehai{788m 799p 113s 33z -666'z}[ドラ \mj{8s}][toitoi-4]

手牌\ref{mj:toitoi-5}はツモ \mj{2p} の場面。トイトイを意識しないと早い巡目でツモ切ってしまうがち。ここは打点のために打 \mj{3p}。4巡目以降は打 \mj{2z} としてトイトイと手なりの道を残す\citep[][p.~58]{zero-teyaku}。

\tehai{8899m 1223p 233z -666'z}[東場　南家　3巡目　ドラ \mj{8s}][toitoi-5]

手牌\ref{mj:toitoi-6}はトイトイとして対子が一つ足りていない。手なりで \mj{2m} を切りたいが、もう一つの対子ができたら5200に近くなるので、まず対子に手をかけない。打 \mj{7p} にすると \mj{8p} の受け入れがなくなるが、\mj{9p} が出やすくなる。\mj{3z} としてもトイトイへの道は残っているが、安全度と重なった時の鳴きやすさが優れているので残したい。よって、ここは打 \mj{7p}。もう一つの対子ができたら \mj{45s} を払っていく\citep[][p.~58]{zero-teyaku}。

\tehai{223m 11799p 45s 3z -77'7z}[ドラ \mj{8s}][toitoi-6]


\subsubsection*{ゴーニー理論}

\textbf{待ちが苦しい8000より両面5200}、というのは福地誠氏が提唱するゴーニー理論\citep{yoteru-toitoi}。例え手牌\ref{mj:toitoi-7}はドラドラ・中・トイトイで満貫だが、待ちが \mj{5p5s} 双碰なのであまり出てこない。\mj{47s} にすると打点が5200になっても待ちが8枚になる。ゴーニー理論によるとここは打 \mj{5s}。

\tehai{111m 05p 056s -999'm 77'7z}[][toitoi-7]

一方、手牌\ref{mj:toitoi-8}もドラ・中・トイトイで満貫。両面にすると打点が2600にすぎないので、ここはトイトイを狙う。打 \mj{6s}。

\tehai{111m 05p 556s -88'8s 77'7z}[][toitoi-8]


\subsubsection*{トイトイの決め打ち}

もっと遠くて、対子も５つ揃っていない階段でトイトイを狙う基準は\citep[][p.~59]{zero-teyaku}、

\begin{itemize}
\item 打点が5200以上、つまり役牌などの役に複合している
\item 対子が４つ、そのうち尖張牌の対子が１つ以下
\end{itemize}

手牌\ref{mj:toitoi-9}は三向聴だが両面が一切なしで実は受け入れがひどく狭い三向聴。真っ直ぐに行けば \mj{7z} ノミで1000にすぎない。よって、ここはトイトイを狙う。そうすると、重なったら鳴きやすいと高い牌を優先する。つまり、尖張牌＞28＞19＞字牌の順番に牌を切る。手牌からポンできる牌を全部鳴いて \mj{46p} を切る。

\tehai{2288m 1146p 18s 5677z}[ドラ \mj{6s}][toitoi-9]

手牌\ref{mj:toitoi-10}もトイトイを狙う。対子の隣の牌は早ければ早く切るほどいい。よって、ここは打 \mj{7m}

\tehai{22799m 1179p 1s 5677z}[ドラ \mj{6s}][toitoi-10]

手牌\ref{mj:toitoi-11}は \mj{9m} をポンしたばかりの場面。トイトイに向かう際で考えるのは自分の河がどう見えている。ポン \mj{9m} した後 \mj{3p} を切ることで相手を混一色を思わせて、\mj{1p} が出やすくなる。よって、ここは打 \mj{17m} ではなく打 \mj{3p}。\mj{9m} ではなく \mj{1p} からポンすると打 \mj{7m} になる。

\tehai{1227m 1139p 677z -99'9m}[ドラ \mj{6s}][toitoi-11]


\subsubsection*{トイトイにしないケース}

手牌\ref{mj:toitoi-12}は対子4つあっても、トイトイや三暗刻や四暗刻にしない。なぜなら、今回のドラは \mj{5m}、しかも自分が持っているのは \mj{0m}。トイトイにして鳴いたら \mj{5m} が重ならない限り手牌の価値は2600に過ぎない。こうなるとタンヤオで \mj{0m} を使い切った方がいい。そして安全度のために \mj{1s} を先に切る\citep[][pp.~128--129]{fkt}。

\tehai{0m 45666p 112246s 44z}[東場　子　3巡目　ドラ \mj{5m}][toitoi-12]


%---------------------------------------------
\subsubsection{七対子}

\subsubsection*{三対子から七対子へ}

七対子は一向聴の時受け入れが9枚しかないので、通常は狙うものではないし、決め打ちもしない。対子が３つ持っていても七対子のルートを残す場面は手牌\ref{mj:chit-1}のような嵌張だらけのバラバラ牌姿。対子が３つなので七対子三向聴で受け入れが21枚。面子手でも \mj{3p5z} どちらを切ると受け入れが \mj{15m8p456s3z} 22枚になる。肝心なのは、面子手で進むと愚形立直・ドラ1で済む可能性が高い。それに対して、七対子にすると最大立直・ツモ・七対子・裏裏の倍満まで見える。よって、ここは対子をほぐさず孤立している \mj{3p} を切る\citep[][pp.~82--83]{zero-teyaku}。

\tehai{1146m 379p 3557s 335z}[東場　北家　3巡目　ドラ \mj{9p}][chit-1]

手牌\ref{mj:chit-1}で次巡で \mj{9p} をツモって\ref{mj:chit-2}になる。ここでも打 \mj{37s} として七対子に決め打ってはいけない。面子手ルートも並行して進めるので、打 \mj{4m} \citep[][p.~83]{zero-teyaku}。

\tehai{1146m 799p 3557s 335z}[東場　北家　4巡目　ドラ \mj{9p}][chit-2]


\subsubsection*{好形ばかりなら七対子は基本的無視}

手牌\ref{mj:chit-3}のような対子ばかりの手、面子手両向聴で七対子一向聴。ですが、対子にフォロー牌があって、面子手にしても聴牌までそんなに遠くない。対子のフォロー牌を外すか、対子をほぐすか、というのはこの小節の課題。

手牌\ref{mj:chit-3}には役牌暗刻があって、仕掛けが効いている。対子が多くて受け入れは多少狭くて苦しいが、まずわざと七対子のために役牌暗刻を外すことはない。役牌暗刻以外の部分には対子と両面対子しかいない場合、七対子を無視して両面を固定する。手牌\ref{mj:chit-4}のように嵌張対子があったらそこからフォロー牌を外して七対子手と面子手の並行をみる。\mj{1p} を切ると4枚の \mj{2p} がロスになるのに対して \mj{7m6p} 6枚で七対子聴牌するので、受け入れは実は増える\citep[][p.~83]{zero-teyaku}。

\tehai{788m 334677p 05s 777z}[東家　5巡目　ドラ \mj{2s}][chit-3]

\tehai{788m 133677p 05s 777z}[東家　5巡目　ドラ \mj{2s}][chit-4]

手牌\ref{mj:chit-5}は二度受けの並び対子があって、序盤なら七対子を無視して打 \mj{4m} とする。中盤以降なら打 \mj{78m} として面子手と対子手を並行させる\citep[][p.~84]{zero-teyaku}。

\tehai{334478m 22388s 777z}[ドラ \mj{1p}][chit-5]

手牌\ref{mj:chit-6}のようにマンズの上が嵌張になると巡目に関わらず打 \mj{7m} として面子手と対子手を並行させる\citep[][p.~84]{zero-teyaku}。

\tehai{334478m 22388s 777z}[ドラ \mj{1p}][chit-6]

手牌\ref{mj:chit-7}は表~\ref{tab:chit-7}の通り \mj{8p6s} のどちらを切っても受け入れが同じ6種17枚。\mj{8p} 切ったら \mj{2233p} の並び対子が両面＋両面にならないので \mj{1p} がロスになるが、\mj{1p} 引いたらタンヤオが崩すのでさほど痛くない。\mj{8p} 切りにより七対子の一向聴も残る。ツモ次第でタンヤオへ移行することもできる\citep[][p.~18]{uzaku-nankiru}。

\tehai{2233556678p5667s}[東1局　東家　5巡目][chit-7]

\begin{table}[ht]
  \centering
  \begin{tabular}{ll}
    切る牌 & 受け入れ \\
    \hline
    \mj{8p} & \mj{2347p57s} 17枚 \\
    \mj{6s} & \mj{123457p} 17枚
  \end{tabular}
  \caption{手牌\ref{mj:chit-7}の手の受け入れ。}
  \label{tab:chit-7}
\end{table}

手牌\ref{mj:chit-8}にあるドラダブ東を鳴きたいだが、もう9巡目、そして \mj{1z} を鳴いても向聴数が進まない。よって、ここは面子手と対子手を並行させる。打 \mj{56p} という決め打ちはしない。そうすると、\mj{3m} と \mj{7p} の比較は \mj{1z} を鳴いた後の形に基づく。表~\ref{tab:chit-8} の通り、打 \mj{7p} の方の受け入れが2枚多い。打 \mj{7p} とすると \mj{258p} の三面張が崩れたけれど、両面対子＋両面対子＋対子の方は二度受けの両面両面＋対子＋対子の方より受け入れが広い\citep[][p.~20]{uzaku-nankiru}。

\tehai{344m 3344567p 55s 11z}[東1局　東家　9巡目　ドラ \mj{1z}][chit-8]

\begin{table}[ht]
  \centering
  \begin{tabular}{ll}
    \multicolumn{2}{l}{\mj{44m 3344567p 55s -1'11z}} \\[3pt]
    打牌 & 受け入れ \\
    \hline
    \mj{4m} & \mj{23458p5s} 17枚 \\
    \mj{3p} & \mj{4m2458p5s} 17枚 \\
    \mj{4p} & \mj{4m2358p5s} 17枚 \\
    \mj{5s} & \mj{4m23458p} 17枚 \\
    \\
    \multicolumn{2}{l}{\mj{344m 334456p 55s -1'11z}} \\[3pt]
    打牌 & 受け入れ \\
    \hline
    \mj{4m} & \mj{25m235p5s} 19枚 \\
    \mj{3p} & \mj{245m25p5s} 19枚 \\
    \mj{6p} & \mj{245m25p5s} 19枚 
  \end{tabular}
  \caption{手牌\ref{mj:chit-8}の \mj{1z} を鳴いた形}
  \label{tab:chit-8}
\end{table}


\subsubsection*{一段目の河に工夫をする}

手牌\ref{mj:chit-9}は手なりなら打 \mj{8m}。面子手と七対子の両天秤にすると打 \mj{59s}。が、ここはまだ一段目なので、七対子一向聴の階段でアガリ率を最大化するため打 \mj{7s} の選択肢も浮上してくる。

\tehai{8899m 33445p 579s 33z}[南家　4巡目　ドラ \mj{1p}][chit-9]

手牌\ref{mj:chit-10}も打 \mj{3m} $\to$ \mj{6m} の手順で \mj{1m} が安全に見える河を作っていく。

\tehai{136m 88p 6699s 22334z}[東家　4巡目　ドラ \mj{1s}][chit-10]

そういっても、二段目に入ると河に中張牌たくさん切られていてやや不自然にになってしまうので、中盤以降は河のことを考えなくてもいい。


\subsubsection*{消極的七対子}

七対子は先手を取られると非常に弱くなる。特に中盤以降、誰かから立直がかかってくるのもおかしくない階段、守備力がある孤立牌を残すのは「消極的七対子」と呼ばれている。手牌\ref{mj:chit-11}からは打 \mj{2z} ではなく、立直を受けた後一向聴を維持するため、中張牌の \mj{6p} にすべき。振聴の \mj{2z} だが、生牌の \mj{6p} と2枚切れの \mj{2z} は重なる確率はそう変わらないし、後手を踏んだ時 \mj{6p} は切りづらい\citep[][p.~85]{zero-teyaku}。

\tehai{499m 67799p 11779s 2z}[東家　10巡目　ドラ \mj{4m}　\mj{2z} は2枚切れで振聴][chit-11]


\subsubsection*{七対子とトイトイの分岐点}

\citet{gendai-toitoi}によると、

\begin{itemize}
  \item トイトイ有利
  \begin{itemize}
    \item 暗刻がある
    \item 役牌や混一色などの役と複合している
    \item 端牌対子で仕掛けやすい
    \item オーラスで僅差でアガリトップ
  \end{itemize}
  \item 七対子有利
  \begin{itemize}
    \item 2枚切れている刻子にならない対子がある
    \item 暗刻がない
    \item 他に役がない
    \item 孤立のドラがある
    \item 対子が鳴きにくい
    \item 余裕のあるトップ目等守備を重視する場合
  \end{itemize}
\end{itemize}

手牌\ref{mj:chit-toi-1}から \mj{8s} を切るのは牌効率通りの思考で、完全一向聴も見える。が、打点が1000にすぎない。アガリトップなどの場面以外なら打 \mj{34p} として七対子やトイトイの可能性を最大限にする。6枚損だが、両向聴なので打点上昇のためにいいトレードオフと思う。ちなみに打 \mj{4m} も七対子とトイトイの道が残っているが、両面対子が消えていて一向聴の形は悪くなる\citep{gendai-toitoi}。

\tehai{33477889m 34p 88s 77z}[][chit-toi-1]

手牌\ref{mj:chit-toi-2}はトイトイ有利の条件と七対子有利の条件が満たされている。七対子ドラ単騎が上がりにくいという点から考えると、トップ目ではないと打 \mj{8s} としてトイトイに向かう\citep{gendai-toitoi}。

\tehai{33388m 1155p 8s 4566z}[ドラ \mj{4z}][chit-toi-2]

\begin{itemize}
  \item トイトイ有利
  \begin{itemize}
    \item 暗刻がある
    \item 役牌対子がある
    \item 尖張牌対子が１つ
  \end{itemize}
  \item 七対子有利
  \begin{itemize}
    \item 孤立のドラがある
  \end{itemize}
\end{itemize}

2暗刻＋4対子の形も、役牌と複合しているか、七対子で即リーできる牌の枚数で判断する。手牌\ref{mj:chit-toi-3}は、\mj{3m8s} とすると \mj{8m15p6z} の他の牌がツモったら七対子聴牌になる。打 \mj{5p} とするとポンテンできるし四暗刻まで見える。今回は \mj{66z} の対子があるので、トイトイに向かう方がいい。

\tehai{33388m 1155p 888s 66z}[ドラ \mj{4z}][chit-toi-3]

%---------------------------------------------------
\subsubsection{三暗刻}

三暗刻聴牌の立直判断に関しては\ref{sec:riichi-snk}を参照。

\subsubsection*{対子３つ＋暗刻１つなら三暗刻を視野に}

手牌\ref{mj:snk-1}は中ぶくれの \mj{2334m} で真ん中の \mj{3m} を引いてしまう結果。\mj{3m} ではなければ \mj{89p} を払ってピンフ一盃口が見えるが、\mj{3m} なら対子系への距離は順子系より近い。ツモ \mj{9m8p9p} でツモり三暗刻、他の部分に対子がもう一つできたら七対子やトイトイも見える。よって、ここは \mj{4m} $\to$ \mj{2m} で切っていく\citep[][p.~73]{zero-teyaku}。

\tehai{2333499m 5678899p}[南家　5巡目　ドラ \mj{3z}][snk-1]

手牌\ref{mj:snk-2}は手なりでいくと \mj{4m8s} が、愚形リーのみがほぼ確定。ここも手なりで進行せず、打 \mj{2m} で三暗刻を狙う。打 \mj{9s} ではないのは安全度のため。手なりで進んでもいいのはドラが2枚以上の時\citep[][p.~74]{zero-teyaku}。

\tehai{244m 12388p 111889s}[南家　5巡目　ドラ \mj{3z}][snk-2]

手牌\ref{mj:snk-3}は節\ref{sec:complex-shape}に触れられたケースと似ていて、両面対子ばかりの手牌は対子のフォロー牌を外して一向聴の広さを瞬間的受け入れより優先する。縦引きで成り立つ三暗刻から考えるとフォロー牌外しも有利。ここは安全度と赤受けを考えて打 \mj{4p} とする\citep[][p.~74]{zero-teyaku}。

\tehai{233m 455p 23677s 777z}[ドラ \mj{3z}][snk-3]


%============================================
\subsection{4ブロック打法}

\textbf{聴牌まで遠い序盤なら辺張・嵌張より役・打点に関連する孤立牌を残すのもあり}\citep[][pp.~158--159]{suphx}。このような5ブロックを4ブロックにするのは「向聴戻し」ともいう。序盤で向聴戻しの有効局面は下記のパターンの一つ\citep{gendai-shantenback}：

\begin{enumerate}
\item 強い形（4連形、中ぶくれ）$\times$2
\item 強い形＋打点増加孤立牌
\item 打点増加孤立牌$\times$２
\item 愚形にフォローがある
\item 愚形が枯れている（2枚切れ）
\item 一向聴に取ると愚形安手確定
\end{enumerate}

手牌\ref{mj:rst-1}は混一色が見えるのでここで外嵌張の \mj{79m} を外す。

\tehai{79m4456789p33445z}[東1局　西家　6巡目　ドラ \mj{1z}][rst-1]

手牌\ref{mj:rst-2}はチャンタや三色が見えるのでここで \mj{77s} を外す。

\tehai{112379m9p77789s44z}[東1局　西家　6巡目　ドラ \mj{1z}][rst-2]

手牌\ref{mj:rst-3}で \mj{46m} には \mj{33m} のフォローがあるので \mj{6m} を切っても \mj{5m} の受けは逃さない。

\tehai{334699m36778p567s}[東1局　西家　6巡目　ドラ \mj{1z}][rst-3]

手牌\ref{mj:rst-4}ではそのまま一向聴をとると外嵌張や辺張待ちになるのでここで \mj{89m} を落とす。

\tehai{1389m224566p4567s}[東1局　西家　6巡目　ドラ \mj{1z}][rst-4]

手牌\ref{mj:rst-5}は \mj{7m} が枯れているなら \mj{89m} を落とす。\mj{7m} まだ4枚あるならそのままで \mj{6p} を切る。

\tehai{23489m11566p4567s}[東1局　西家　6巡目　ドラ \mj{1z}][rst-5]

手牌\ref{mj:kanchan-1a}にはタンヤオがつかないため、辺張を落としても打点はよほどよくならなさそう。聴牌のときの最終形の想定すれば、辺 \mj{7p} を避けるもっといい待ちを作る\citep[][pp. 48--55]{tetsusenryaku}。

\tehai{4678m 3589p 11067 -7s}[東1局　西家　3巡目　ドラ \mj{2m}][kanchan-1a]

ドラ3で打点はもう十分である場合は、素直に聴牌に向かうべき。ですが、ソウズの部分が \mj{114067s} になったら伸びやすいので辺張を外すのはおすすめ。

\tehai{4678m 3589p 11067 -7s}[東1局　西家　3巡目　ドラ \mj{1s}]

手牌\ref{mj:block-11}に愚形がばかり。そのままで立直すれば大体2600に過ぎない。点数が欲しい場面ならチャンタを見込んで、\mj{5s} の対子を落とし、\mj{2z} 重なりや \mj{4z} 槓で打点を最大限にする\citep[][pp. 74--79]{tetsusenryaku}。

\tehai{123 89m 79p 155s 244 -4z}[南4局　南家　1巡目　ドラ \mj{1s}][block-11]

Suphxは\ref{mj:block-9a}から \mj{13m} を落とした\citep[][pp.~158--159]{suphx}。これらのブロックをキープしたら、仕掛けるとほぼ2飜、門前も結構厳しめ。\mj{0p} のくっつきにより1飜の打点上昇や、\mj{48s7z} のくっつきにより混一色への転移は十分見込める

\tehai{137m 0p 2248s 22667z -9m}[東2局　北家　1巡目　ドラ \mj{8m}][block-9a]

\textbf{打点上昇を諦めて愚形確定かつ低打点の一向聴を取ることはない}\citep[][pp.~160--161]{suphx}。手牌\ref{mj:block-10a}は一向聴だがそのままで聴牌を取るとドラ \mj{5m} や \mj{0s} が出ていくのは確定。Suphxはここで \mj{5p} を落とした。ただし、\mj{0s} が \mj{5s} となる場合、向聴を戻しても打点が必ず上昇するわけではないため、向聴を戻さず立直に寄る。

\tehai{588m 1135p 130s 777z -2s}[南2局　北家　5巡目　ドラ \mj{5m}][block-10a]

麻雀は4面子1雀頭を作るゲーム。打点を上げるためや、よりいい形を求めるために瞬間のブロック不足は状況により許容されるが、三向聴・四向聴の手牌の受けがどんなに広いでも聴牌までのツモ巡数は必ず増えていく。前述の4ブロックにする打法はいずれも打点・手役に関連するので、いつもドラを貪って愚形を拒否するわけではない。

%============================================
\subsection{一向聴(1)：種類}

麻雀は14枚の牌で「4面子＋1頭」をできるゲームなので、聴牌の時は「頭なし4面子」（単騎・延べ単）と「3面子＋1頭＋1搭子」（辺張・嵌張・両面）もしくは多面張。一向聴は、前述の形をもう一回崩して、その種類は下記の通り\citep{tomo-iishanten}。

\begin{table}[ht]
  \centering
  \begin{tabular}{lllll}
    & 面子 & 対子 & 他の搭子 & 浮き牌 \\
    \hline
    余剰牌一向聴 & 2 & 1 & 2 & 1 \\
    完全一向聴 & 2 & 1 & 2 & 1 \\
    ヘッドレス一向聴 & 3 & 0 & 2 & 0 \\
    くっつき一向聴 & 3 & 1 & 0 & 2
  \end{tabular}
  \caption{一向聴の種類}
\end{table}


\subsubsection*{余剰牌一向聴}

余剰牌がある一向聴は、面子が二つ、搭子（対子含み）が二つがある上、あと一つの孤立牌が浮いている状態を指す。その余剰牌の機能は下記の通り三つある：

\begin{enumerate}
  \item 安全牌
  \item 良形変化
  \item 打点上昇
\end{enumerate}

\tehai{12389m12366p78s-?}[][ist-1]

手牌\ref{mj:ist-1}の \mj{?} が \mj{4z} の場合、誰にも不要な字牌を抱えて守備へ寄っている。\mj{4s} の場合、\mj{35s} 引いたら良形ができ \mj{89m} の辺張を払う。\mj{2s} の場合、\mj{13s} 引いたら下の三色が見える。余剰牌が出ているため、受け入れが狭い。手牌\ref{mj:ist-1}の受け入れは \mj{7m69s} 最大3種12枚。\mj{89m} が \mj{78m} になっても受け入れが4種16枚に過ぎない。

手牌が良形確定なら余剰牌を安牌としてを抱えていく。愚形があるなら自分の手で優先して良形変化への余剰牌を持っていく。


\subsubsection*{完全一向聴} 

\tehai{123789m23566p88s}[][ist-2]

\tehai{123788m23456p88s}[][ist-3]

完全一向聴は\ref{mj:ist-2}のように、その浮き牌がある両面搭子に重なってできる両面対子がある形。その受け入れが、ある両面搭子が両面対子のともに二度受けになるという最低の場合でも5種16枚。両面搭子が順子に繋がって三面張ができる場合、例え\ref{mj:ist-3}のように、受け入れが \mj{689m147p8s} 7種23枚となる。

\tehai{22m566788p34678s}[][ist-3-1]

\tehai{688m22p567p34678s}[][ist-3-2]

手牌\ref{mj:ist-3-1}が含んでいる \mj{566788p} の形があれば完全一向聴になる。受け入れは \mj{2478p25s} 6種19枚。この形は \mj{567p}＋\mj{688p} で分析できるが、19枚の受け入れになる「順子＋嵌張対子」のはこの形だけ。例え\ref{mj:ist-3-2}のように \mj{688p} が \mj{688m} になると受け入れが \mj{78m2p25s} 5種16枚となる。

\tehai{12378m223456p88s}[][ist-3-3]

手牌\ref{mj:ist-3-3}のように、ピンズの部分が \mj{223456p} か \mj{233456p} か、両面対子が順子に繋がった場合受け入れが7種23枚となる。

\subsubsection*{ヘッドレス一向聴} 

\tehai{123789m23567p78s}[][ist-4]

\tehai{111789m23567p78s}[][ist-5]

\tehai{111789m34567p78s}[][ist-6]

ヘッドレス一向聴は\ref{mj:ist-4}のように全ての搭子には対子がない形。\mj{23p78s} のどちらが先に重なったら良形聴牌になるが、\mj{14p69s} 先に引いたら単騎になってしまうリスクがある。従って、手牌\ref{mj:ist-4}の受け入れが \mj{1234p6789s} 8種28枚となるが、良形聴牌できるツモ牌はただ4種12枚。

暗刻が含まれているヘッドレス一向聴は必ず良形聴牌になる。例え\ref{mj:ist-5}のように、\mj{14p69s} 引いても \mj{1m} 切って残りの両面が十分に働く。よって、\ref{mj:ist-4}と\ref{mj:ist-5}の受け入れが同じ8種28枚であるが、良形聴牌率によると後の方が圧倒的に有利。そして\ref{mj:ist-6}のようにある両面が順子につながって三面張ができたら受け入れが一気に11種37枚になるので、非常に強い形である。


\subsubsection*{くっつき一向聴}

\tehai{12355789m3789p3s}[][ist-7]

\tehai{12355789m1789p1s}[][ist-8]

くっつき一向聴は\ref{mj:ist-7}のように、二つの浮き牌と頭に関連する牌が引いたら聴牌できる形。その受け入れが11種40枚（\mj{5m} + \mj{12345p} + \mj{12345s}）あり、一番聴牌になりやすい形であるが、良形聴牌できる牌は \mj{24p24s} 4種16枚だけ。さらに\ref{mj:ist-8}のように浮き牌が端牌になる場合、受け入れが7種24枚になり、有効牌引いても必ず愚形聴牌になる。

%============================================
\subsection{一向聴(2)：選択}

ここまでの話は全ての搭子に愚形がないという前提の上で成立する。愚形があったら受け入れがさらに減って、不十分と思える形。聴牌になりにくくて、聴牌になっても高い確率で愚形聴牌になるので、\textbf{不十分の一向聴の場合で考えるのは良形変化と向聴戻し}\citep{manboo-iishanten}。

\tehai{122357m5678p2245s}[][ist-9]

手牌\ref{mj:ist-9}では \mj{2m58p} どちら切っても受け入れが \mj{6m36s} 3種12枚になるが、\mj{58p} を残して \mj{4679p} を引いたら良形になるのでここで \mj{2m} を切る。

\tehai{1223467p2345778s}[][ist-10]

手牌\ref{mj:ist-10}では \mj{258s} どちら切っても受け入れが \mj{358p} 3種11枚になる上、一枚の \mj{3p} が既に手牌にあるので、巡目がまだ深くないなら打 \mj{2p} で両向聴に戻ってよりいい一向聴を取るのは得。

\tehai{4578999m45p11145s}[][ist-11]

手牌\ref{mj:ist-11}では \mj{45m}・\mj{45p}・\mj{45s} の三択。どちらを切っても受け入れに差が出ないけれど、\mj{45m} を落とした場合 \mj{6m} を引いたら暗刻ありヘッドなしの広い一向聴になるのでそっちの方がより有利。

\tehai{4668m34445p23777s}[][ist-14]

手牌\ref{mj:ist-14}では、打 \mj{6m} の受け入れ（\mj{57m14s} 4種16枚）が打 \mj{8m} の受け入れ（\mj{56m4p14s} 5種15枚）により1枚多いけれども、打 \mj{8m} の後に \mj{2356p} 引いて \mj{4m} 切りで完全一向聴ができるので、受け入れの枚数に差が大きくない場合良形変化を優先させる。

\tehai{1234m 5678p 122234s}[][ist-28]

\begin{table}[ht]
  \centering
  \begin{tabular}{ll}
    打牌 & 受け入れ \\
    \hline
    \mj{1s} & \sfrac{14}{46}：\mj{123456m3456789p3s} \\
    \mj{1m} & \sfrac{7}{20}：\mj{3456789p12345s}
  \end{tabular}
  \caption{手牌\ref{mj:ist-28}の受け入れ}
  \label{tab:ist-28}
\end{table}

手牌\ref{mj:ist-28}は表~\ref{tab:ist-28} の通り打 \mj{1s} の受け入れが最も広いが、ピンフや一盃口に向かうなら打 \mj{1s} にする\footnote{\mjfn{2s} がドラなら打 \mjfn{1s} にして一番早い速度で聴牌をとる。}。\mj{122233s} の形は \mj{123s} + \mj{223s} に見えるのでピンフ・一盃口になりやすい形\citep[][p.~14]{uzaku-nankiru}。

\tehai{678m 123306p 12378s}[][ist-29]

\begin{table}[ht]
  \centering
  \begin{tabular}{ll}
    打牌 & 受け入れ \\
    \hline
    \mj{3p} & \sfrac{8}{28}：\mj{4567p6789s} \\
    \mj{12p} & \sfrac{4}{16}：\mj{47p69s}
  \end{tabular}
  \caption{手牌\ref{mj:ist-29}の受け入れ}
  \label{tab:ist-29}
\end{table}

手牌\ref{mj:ist-29}は表~\ref{tab:ist-29} の通り打 \mj{3p} とするヘッドレス一向聴の受け入れが最も広いが、半分以上の確率で望ましくない単騎待ちになる。\mj{3p} を対子にして \mj{12p} を落としても \mj{47p69s} の受け入れも消えない上、ツモ \mj{6p} で \mj{0p} が出ていくのを避けるためにここでヘッドレス一向聴にしない\citep[][p.~16]{uzaku-nankiru}。


\subsubsection*{強い待ちを残す}

\tehai{22356799p234556s}[][ist-12]

手牌\ref{mj:ist-12}では \mj{2p} と \mj{5s} の二択。どちらの受け入れも同じだけど、\mj{556s} に \mj{5s} の受け入れを残した場合 \mj{5s9p} の双碰受けに引いたら三面張が崩れてしまうのでここで打 \mj{5s} として三面張を固定する。

\tehai{223m22789p445s -7'77z}[][ist-31]

手牌\ref{mj:ist-31}では \mj{2m} と \mj{4s} の二択。\mj{2m} の方が仕掛けやすい（聴牌できる確率がより高い）が、\mj{14m} 待ちの和了できる率がより高い。どっちにせよ、聴牌率は常に上がり率より高いので、上がり率を重視してここで \mj{2m} を切る。


\subsubsection*{役による打牌判断}

選択肢同士の受け入れに差が出ない場合、役で判断する。

\tehai{67m45688p3345556s}[][ist-13]

\begin{table}[ht]
  \centering
  \begin{tabular}{ll}
    打牌 & 受け入れ \\
    \hline
    \mj{3s} & \sfrac{7}{22}：\mj{58m8p2457s} \\
    \mj{6s} & \sfrac{7}{20}：\mj{58m8p2345s} \\
    \mj{5s} & \sfrac{6}{19}：\mj{58m8p347s}
  \end{tabular}
  \caption{手牌\ref{mj:ist-13}の受け入れ}
\end{table}

手牌\ref{mj:ist-13}の場合、\mj{5s} を切って \mj{334455s} の一盃口になる可能性があるが打 \mj{3s} の受け入れと良形聴牌率で比べると不利なのでここで一盃口を見切って打 \mj{3s} とする。

\tehai{12355m1334556p23s}[][ist-15]

手牌\ref{mj:ist-15}では打 \mj{1p} と打 \mj{3p} の受け入れが同じだが、打 \mj{1p} の方に一盃口チャンスがあり、一方打 \mj{3p} の方に三色とピンフの可能性がある。よって、ここで打 \mj{3p} とする。

\tehai{11155m1334556p23s}[][ist-16]

手牌\ref{mj:ist-16}は\ref{mj:ist-15}の変化形、三色もピンフもつかない。この状況なら \mj{1p} を切って、ツモ \mj{5p} で完全一向聴を期待する。

\tehai{1222m112378p5567s}[][ist-19]

手牌\ref{mj:ist-19}に打 \mj{1m} に受け入れが20枚、打 \mj{1p} に21枚、打 \mj{5s} にも21枚ある。打 \mj{1m} とするとピンフにならない上 \mj{13m} のような良い待ちを失うのでここで論外。亜両面は内側の方に良形変化の枚数が多いのでここで \mj{1p} を切る。

%\tehai{12357889m067p556s}[][ist-32]
%
%\begin{table}[ht]
%  \centering
%  \begin{tabular}{ll}
%    打牌 & 受け入れ \\
%    \hline
%    \mj{8m} & \sfrac{10}{34}：\mj{34567m45678s} \\
%    \mj{9m} & \sfrac{5}{16}：\mj{68m457s}
%  \end{tabular}
%  \caption{手牌\ref{mj:ist-32}の受け入れ}
%  \label{tab:ist-32}
%\end{table}
%
%手牌\ref{mj:ist-32}は一通（\mj{5m} + \mj{789m}、\mj{8m} が不要）と三色（\mj{57m} + \mj{88m}、\mj{9m5s} が不要）の選択\footnote{そう言ってもどちらの役はまだ確定されていない。}。\mj{8m} を切ったら \mj{5m6s} のくっつき一向聴なので表~\ref{tab:ist-32} 打 \mj{9m} の方より受け入れが広い。よって、\citet[p.~18]{uzaku-nankiru}は打 \mj{8m} にする。
%
%\begin{table}[ht]
%  \centering
%  \begin{tabular}{lll}
%     & 価値 & 確率 \\
%    \hline
%    打 \mj{8m} & 2600 & 50\%（ツモ \mj{37m48s}） \\
%     & 3900 & \minibox{38.2\%（ツモ \mj{5m57s} \\ + ツモ \mj{46m} $\times$ 50\%）}\\
%     & 7700 & 11.8\%（ツモ \mj{46m} $\times$ 50\%）\\
%    打 \mj{9m} & 2600 & 50\%（ツモ \mj{8m45s}）\\
%     & 3900 & 12.5\%（ツモ \mj{6m} $\times$ 50\%）\\
%     & 8000 & \minibox{37.5\%（ツモ \mj{7s} \\ + ツモ \mj{6m} $\times$ 50\%）}
%  \end{tabular}
%  \caption{手牌\ref{mj:ist-32}の価値の期待値}
%  \label{tab:ist-32-ev}
%\end{table}
%
%実際それぞれのツモとアガリ形の価値の期待値を求めると表~\ref{tab:ist-32-ev} の通り三色の方の期待値が高いのがわかる。問題になるのは受け入れを34枚を16枚にすることによるツモ巡数の増加による局収支はどれくらい下がるか。役が確定されていない上、\mj{8m} を切って両面聴牌になっても待ちが打 \mj{9m} の方に明らかに優れていないので、やはり場面によるか\footnote{東1局自分西家7巡目平場の場面なら打 \mj{8m} にする、それより点数が欲しい場面なら打 \mj{9m}、と。}。


\subsubsection*{愚形含み完全一向聴}

普通には安牌を持つためにも良形完全一向聴を余剰牌一向聴にはしない\footnote{受け入れが枯れている場合なら安牌を残す。}\citep{gendai-iishanten-choice-1}。単独搭子が愚形である場合、選択肢が下記のように三つある。

\begin{enumerate}
\item 複合搭子を両面にして余剰牌一向聴を取る
\item 浮き牌を切って愚形含み完全一向聴を取る
\item 愚形搭子を払って両向聴に戻る
\end{enumerate}

\tehai{2347m5589p123788s}[][ist-201]

手牌\ref{mj:ist-201}は \mj{7m} を切ると「愚形含み」完全一向聴になる。\mj{8s} 切りで余剰牌一向聴になる。\mj{7m} がドラそばの場合、\mj{8s} 切りでドラ絡みの嵌張・両面を期待する。そうでなければ \mj{7m} 切りで素直に愚形含み完全一向聴を取る方がいい。

\tehai{2347m4699p123788s}[][ist-21]

手牌\ref{mj:ist-21}なら愚形が内嵌張になって両面変化ができやすいので \mj{7m} がよほど重要ではないのでここで切る。

\tehai{2345m5589p123788s}[][ist-22]

手牌\ref{mj:ist-22}なら孤立牌が順子についているので良形ができやすい。よって、ここで \mj{8s} を切って6ブロックにする。

\tehai{2345m6689p123788s}[][ist-23]

手牌\ref{mj:ist-23}でさらに頭が \mj{6p} になって \mj{7p} のフォローができるので、直接に \mj{89p} を落としても \mj{7p} がロースにならないので向聴戻ししてもいい。

\tehai{2345m4699p123788s}[][ist-24]

四連形と内嵌張の組み合わせ。極稀の２３４三色のために打 \mj{5m} にする。打 \mj{8s} で6ブロックにしても良いか \mj{37p} 引きで完全一向聴にならないのは損。

\tehai{234m45589p123788s}[][ist-25]

手牌\ref{mj:ist-25}に両面対子が二つあるが、愚形が対子になる変化が期待できるので両面対子を対子にしない。\mj{0p} の受け入れと \mj{69s} という強い待ちを考えたら \mj{8s} を切る。

\tehai{33455m 668p 345667s}[][ist-26]

手牌\ref{mj:ist-26}は打 \mj{6s} ツモ \mj{7p} でヘッドレス一向聴になる可能性があるが、打 \mj{6s} の受け入れが打 \mj{8p} より4枚損。親なら、もしくは巡目が既に深いなら素直に打 \mj{8p} で愚形含み完全一向聴にする\citep[][p. 10]{uzaku-nankiru}。

\tehai{135m 12399p 123667s}[][ist-27]

手牌\ref{mj:ist-27}は打 \mj{5m} と打 \mj{6s} の受け入れは同じ16枚。今回123の三色があるので、打 \mj{5m} により三色を確定しても受け入れが減らないのでここで打 \mj{5m} にする。さらに先切り \mj{5m} によって \mj{2m} が出やすくなる\citep[][p. 12]{uzaku-nankiru}。

\subsubsection*{愚形含みヘッドレス一向聴} \tehai{234888m57p112367s}[][ist-17]

手牌\ref{mj:ist-17}で打 \mj{1s} で広いヘッドレス一向聴ができるが、ソーズの \mj{5678s} を引いたら必ずカン \mj{6p} 待ちになる。よって、巡目がまだ早いなら愚形の崩してよりいい形を待つのは有利。ただし、聴牌を遅らせる余裕がないなら素直に \mj{1s} を切るべき。

\tehai{234888m57p112389s}[][ist-18]

手牌\ref{mj:ist-18}には良形聴牌ができないので、素直に \mj{1s} を切って聴牌を目指す。その時立直するか黙聴するか考えればいい。

%============================================
\subsection{手牌のスリム化}

速度・守備力が重視されるトップ目、役なし愚形濃厚などの場面なら、手牌の受け入れを少し減らしてもあまりロスにはならない。こういった手狭に構えるのは手牌のスリム化ともいう\citep[][p. 29]{suphx}。

点棒状況を無視して、単純に巡目と向聴の数と質で判断しても、\textbf{二段目に入ってまだ両向聴のままなら安牌を残す}\citep[][pp.~106-107]{fkt}。6巡目の時点で、一向聴ならアガれる率は25\%、両向聴だったら13\%。手牌\ref{mj:slim}のような愚形だらけの牌姿ならなおさら。ここは打 \mj{4s}。

\tehai{2368m 24688p 3334s 3z}[東場　子　8巡目　ドラ \mj{1p}][slim]

立体図~\ref{rittai:slima} は \mj{3s} をツモって、受け入れを最大限にするか安牌の \mj{3z} を切るかの問題。ここは \mj{3s} を残さない理由は、

\begin{itemize}
\item \mj{35s} の縦引きでリーのみになる
\item 下家が立直してくると \mj{3s} は切りづらい
\item チー \mj{3m} や立直を宣言する時 \mj{3s} を打つと \mj{2s} が出でづらい
\end{itemize}

\begin{rittai}[h]
  \centering
  \captionsetup[table]{labelsep=space}
  
  \drawrittai{
    jicha/seat = 東,
    jicha/hand = 123m 12356p 13355s 3z,
    jicha/river = 9m  8s 4z 2z 5^z 2p,
    jicha/riverb = 1p,
    jicha/point = 27600,
    %
    toimen/river = 1m 8p 6z 9^p 4s 2s,
    toimen/riverb = 6z,
    toimen/point = 33100,
    %
    kami/river = 1z 9s 9p 7p 4p 9s,
    kami/riverb = 6s,
    kami/point = 16000,
    %
    shimo/river = 9m 3z 7z 8p 4z 6p,
    shimo/riverb = 1s,
    shimo/point = 23300,
    %
    kyoku = 南1局,
    dora = 6z
  }
  
  \caption{}
  \label{rittai:slima}
\end{rittai}


\subsubsection*{序盤のスリム化}

\textbf{トップ目なら鳴ける手役の重要度は高まる}。フラットの状況なら一般の牌効率を従って字牌を先に処理するだろう。立体図~\ref{rittai:slim-1} は南2局のトップ目なので、三色と役牌の重要度と速度は愚形リーのみより高いので、Suphxはここで \mj{2s} を切った\citep[][pp.~202--203]{suphx}。

\begin{rittai}[h]
  \centering
  \captionsetup[table]{labelsep=space}
  
  \drawrittai{
    jicha/seat = 南,
    jicha/hand = 1205678m 13p 223s 2 -6z,
    jicha/river = 3z,
    jicha/point = 37500,
    %
    toimen/river = 1z,
    toimen/point = 9800,
    %
    kami/river = 9m 2z,
    kami/point = 23800,
    %
    shimo/river = 9m,
    shimo/point = 28900,
    %
    kyoku = 南2局,
    dora = 2m
  }
  
  \caption{}
  \label{rittai:slim-1}
\end{rittai}

\textbf{遠いリーのみは追わない}\citep[][pp.~204--205]{suphx}。立体図~\ref{rittai:slim-2} の手牌は数牌の横引きが狭いので、立直までは遠い上、三色・チャンタ・役牌など崩れてしまう可能性も高い。序盤にも関わらずSuphxは \mj{4m} を切った。

\begin{rittai}[h]
  \centering
  \captionsetup[table]{labelsep=space}
  
  \drawrittai{
    jicha/seat = 東,
    jicha/hand = 49m 788999p 1239s 1 -3z,
    jicha/river = 46^5z,
    jicha/point = 41000,
    %
    toimen/river = 6z 1m 3z,
    toimen/point = 17800,
    %
    kami/river = 47z 1^p,
    kami/point = 15000,
    %
    shimo/river = 6z 1m 9^p,
    shimo/point = 26200,
    %
    kyoku = 東4局,
    dora = 3z
  }
  
  \caption{}
  \label{rittai:slim-2}
\end{rittai}

\textbf{明確に高打点かつ早そうな相手に守備を意識すること}\citep[][pp.~206--207]{suphx}。面子手に対する「捨て牌が早そう」はなかなか当たらない。特に序盤が七対子・混一色などの手順により、そうでないパターンが多いだから。早そうに見えたとしても判断に影響するほどならない。

立体図~\ref{rittai:slim-3} で満貫聴牌の手牌ならなおさら。けれど相手が役牌を暗槓して強搭子（中張牌の搭子とかポン材として優れているの \mj{1m} 対子とか）を落とすとなればほぼ早いのは確定される。ここでSuphxは \mj{8m} の受け入れ\footnote{\mjfn{455679m} は両面嵌張。節\ref{subsec:rmkc}を参照。}を捨てて \mj{9m} を落とした。

\begin{rittai}[h]
  \centering
  \captionsetup[table]{labelsep=space}
  
  \drawrittai{
    jicha/seat = 南,
    jicha/hand = 45579m 2206p 345s 4z -6m,
    jicha/river = 57z 8p 3z,
    jicha/point = 25000,
    %
    toimen/hand = XXXXXXXXXX -X66Xz,
    toimen/river = 9^6s 11m,
    toimen/point = 30700,
    %
    kami/river = 4z 1p 3s 5p 4^s,
    kami/point = 18300,
    %
    shimo/hand = XXXXXXXXXX -055's,
    shimo/river = 1s 43z 9s 8^m,
    shimo/point = 25000,
    %
    kyoku = 東3局,
    dora = 3p9m,
    honba = 1,
    kyoutaku = 1
  }
  
  \caption{北家は \mj{6s} の次巡で \mj{6z} カンツモ切り \mj{5s}}
  \label{rittai:slim-3}
\end{rittai}


\subsubsection*{両面対子を崩すケース}

\textbf{両面対子でも手牌により両面が崩れやすいケースは要注意}\citep[][pp.~208--209]{suphx}。立体図~\ref{rittai:slim-4} は恵まれた手牌だが、Suphxは \mj{5s} を先に切った。頭がないため同じ \mj{25m25p} 4種16枚の一向聴だが、\mj{566s} の形を保って \mj{47s} をツモれば暗刻あり頭なしの超広い一向聴になる。

だとしても、ツモ \mj{34m34p} なら \mj{47s} 待ちはより弱い待ちなので \mj{566s} を頭にしなければならない。ツモ \mj{25m25p} なら頭がないため \mj{5s} が出ていかなければならない。つまり、直接の受け入れにも限らず、有効変化 \mj{34m34p47s} の中でも多くの場合は \mj{5s} はいらない。場に一枚切れの \mj{3z} に比べると先に切るがいい。

\begin{rittai}[h]
  \centering
  
  \drawrittai{
    jicha/seat = 西,
    jicha/hand = 34789m 4p 111566s 3z,
    jicha/river = 675z,
    jicha/point = 25000,
    %
    toimen/river = 32z 8p 2^z,
    toimen/point = 25000,
    %
    kami/river = 675z 1m,
    kami/point = 25000,
    %
    shimo/river = 6z 9p 4z,
    shimo/point = 25000,
    %
    kyoku = 東1局,
    dora = 9p,
  }
  
  \caption{}
  \label{rittai:slim-4}
\end{rittai}


\subsubsection*{完全一向聴に構えないケース}

完全一向聴は普通の両面＋両面より4枚の受け入れがあるので多くの人は迷わず完全一向聴に取るのだろう。実際20枚受け入れの完全一向聴の聴牌までの平均巡目は両面＋両面より0.35巡だけ減っている\footnote{作者のシミュレーションによる数値。20枚の受け入れの一向聴は平均的に聴牌までが3.84巡かかる。一方16枚の受け入れは4.19巡。}上、手役・打点・守備力・点棒状況などの要素を加えたら「完全一向聴に取るべき」ということを過信すべきではないのがわかる。\citet[]{yuusee-complete-ist}によると、下記の条件が２つ当てはまれば完全一向聴を外す：

\begin{itemize}
\item 双碰受けが2枚以上切れている
\item 切る牌が34567
\item ピンフや三色などの手役がついている
\item 中盤
\end{itemize}

\tehai{123367m 225p 3788s -8m}[東1局　東家　4巡目　ドラ \mj{4m}][slim-5]

\textbf{聴牌まで遠いなら完全一向聴に構えない}\citep[][pp.~214--215]{suphx}。手牌\ref{mj:slim-5}にはいくつの浮き牌があって、それらの昨日は下記の通り：

\begin{itemize}
\item \mj{3m}：一盃口やドラの受け
\item \mj{5p}：\mj{3467p}の受け
\item \mj{3s}：\mj{1245s}の受け、最終形の待ちとして \mj{5p} より優れている
\item \mj{8s}：一向聴の時 \mj{788s} ＋対子の完全一向聴の構え
\end{itemize}

古来の完全一向聴主義なら \mj{788s} の両面搭子を保って \mj{5p} を切ることをよしとする。\mj{5p} を切ると \mj{3467p} の受けがなくなる上、中盤で完全一向聴をとるとすると安牌がない可能性が高くて、暗刻をツモってリーのみになる可能性もある。そのためSuphxは \mj{8s} を外した。

\tehai{223789mm 99p 12356s 4z}[南家　8巡目　ドラ \mj{3m}][ist-20]

手牌\ref{mj:ist-20}は \mj{4z} を切ると完全一向聴になるが、\mj{2m} はドラ表示牌として既に1枚減っていて、点数が欲しい場況ならリーのみを避けるためここで打 \mj{2m} とする\citep[][p. 19]{ishibashi-black}。


%============================================
\subsection{クソ配牌対処法}

三向聴の配牌は一番多い。次いでは四向聴。この2種類の配牌の合計は75\%台になるといいつつも、ただ1局で四向聴以上の配牌しかこない超不運に突き込むときもある。三向聴の配牌のアガリ率は20\%くらいで、五向聴なら15\%くらいで、つまり7回の中に1回しかアガれない。そのほかの6回のために守備力を落とさずに、アガった1回の手牌の価値を高まるのはクソ悪配牌の進行方針。具体に言うと、

\begin{itemize}
\item 字牌は切らずにも鳴かない
\item 両飜以上の役を狙いつつ仕掛けはできるだけしない
\end{itemize}

手牌\ref{mj:kuso-1}は四向聴の手牌。ここで考えるのは手なりで孤立の字牌を切っていくのではなく、

\begin{itemize}
\item 123の三色がある
\item 字牌が重なったら七対子が視野に入る
\item 字牌が重ならなくても守備に構える
\item 孤立の数牌の機能は（ア）縦引きと（イ）両面を作ること
\end{itemize}

特に最後のポイントで考えると \mj{8s} は \mj{6p} より先に出て行ってもいい。両面ができる牌は \mj{7s} だけから。123の三色から考えると \mj{1p} も要らないが、七対子の道を残すため他の部分から先に処理する\citep{ooi-badhaipai}。

\tehai{12m 1126p 13558s 234z}[南1局　東家　2巡目　ドラ \mj{6z}][kuso-1]

そして次巡で \mj{5p} を引いた。面子手三向聴になったら七対子を無視してもいい。そうなると、\mj{2z} の縦引きを見つつ、生牌のオタ風の \mj{4z} を切る\citep{ooi-badhaipai}。

\tehai{12m 1126p 13558s 234z}[南1局　東家　3巡目　ドラ \mj{6z}　\mj{3z} 1枚切れ][kuso-1a]

立体図~\ref{rittai:kuso-2} はノーテンでラスる点棒状況で、どうかしても最低限としては聴牌取りたい場面\citep{haneru-badhaipai}。こんな状況で5向聴の配牌をもらった。配牌オリせずにまず考えるのは、

\begin{itemize}
\item 789の三色 $\to$ \mj{9m} が必要
\item チャンタ $\to$ \mj{1m2p4z7z} が必要
\item 役牌 \mj{7z}
\end{itemize}

こう見ると \mj{4m6p} は不要。そして789の三色に向かうとき、もう一つ搭子を作るため \mj{6p} は \mj{89p} と被ったせいで \mj{4m} の先に出てもいい。

\begin{rittai}[h]
  \centering
  
  \drawrittai{
    jicha/seat = 南,
    jicha/hand = 149m 2689p 1179s 45z,
    jicha/point = 22000,
    %
    toimen/point = 20000,
    %
    kami/point = 21800,
    %
    shimo/point = 36200,
    %
    kyoku = 南4局,
    dora = 3z,
  }
  
  \caption{}
  \label{rittai:kuso-2}
\end{rittai}

立体図~\ref{rittai:kuso-3} もオーラスで最下位と微差の点数状況。ここで5向聴の手牌をもらった。この場面で注目したいのは：

\begin{itemize}
\item 対面や下家が上がったら自分はラスらない
\item 流局で自分がノーテンでもゲームはまだ進む
\end{itemize}

つまり、この局面はクソ悪配牌を無理やりアガリに向かわせるべきではない。\mj{7z} 対子あってもポンしないし、他の部分も鳴かない。対面や下家が立直してくると \mj{7z} を連打してオリる。対面や下家が仕掛け始めると2副露までアシストする\citep{shibukawa-haipai}。

\begin{rittai}[h]
  \centering
  
  \drawrittai{
    jicha/seat = 南,
    jicha/hand = 1359m 59p 1126s 1277z,
    jicha/point = 20100,
    %
    toimen/point = 32000,
    %
    kami/point = 19900,
    kami/river = 4z,
    %
    shimo/point = 26100,
    %
    kyoku = 南4局,
    dora = 9s,
  }
  
  \caption{}
  \label{rittai:kuso-3}
\end{rittai}


%============================================
\subsection{絞り}

\renewcommand{\renderOption}{1}

「絞り」は、自分にとって不要だが、他家に鳴かれる可能性のある牌を切らないことを指す\citep{gendai-shibori}。絞り自体は、不要牌を自分の手に残すせいで手牌が狭くなる上、絞られた相手の進行を妨害する結果、残りの二人の相手を有利にしてしまう。相手に鳴かれても、その相手は平均打点が下がることになって、残りの二人から立直がきたらめくり合わせること（横移動）も期待できる。つまり、\textbf{多くの場合は絞りは不要。自分の都合で手牌を進んでいい}。

絞りが有効な局面は下記にように\multicitep{gendai-shibori; zero-teyaku, p.~18}、

\begin{itemize}
\item 仕掛けている相手が大物手で、自分の手牌が見合わないとき
\item トップ目等、親の連荘や順位争いしている相手の上がりだけを阻止するとき
\end{itemize}

たとえば、誰かがドラ么九牌をポンして、捨て牌と副露の部分から混一色や役牌バックと推測できる場面なら、役牌を絞ってもいい。たとえば、自分が南場のトップ目で、下家の親が仕掛けている場合、親の連荘を阻止するため絞って残りの二人のアガリを期待する。

立体図~\ref{rittai:shiboru-1} はドラポンが入った親の上家がいて、\mj{1z} をツモった瞬間。自分の手牌はまだアガリまで遠くて、ドラや役も一切ない。そして親の副露に役が見つからない。捨て牌は色の偏りがあって、染め手の可能性が高い。トイトイや役牌バックも否定されていない。\mj{1z} を切ったらロンと言われることはほぼないが、鳴かされたら親の打点を増やして、4000オールや6000オールになると自分のトップ率が下がる。よって、ここは \mj{1z} を絞って打 \mj{6m} とする。

\begin{rittai}[h]
  \centering  \drawrittai{
    jicha/seat = 南,
    jicha/hand = 36m 4578p 2688s 233z -1z,
    jicha/river = XXXX,
    %
    toimen/river = XXXXX,
    %
    kami/hand = XXXXXXXXXX -111'p,
    kami/river = 4z 1s 2751m,
    %
    shimo/river = XXXXX,
    %
    dora = 5s,
  }
  \caption{見合わない手なら、鳴かれる可能性が高い大物手の相手や親に役牌を軽易に切るべきではない。}
  \label{rittai:shiboru-1}
\end{rittai}


\subsubsection*{鳴かせるための切り遅れ}

節~\ref{subsub:yakuhai} にも言及されて、役牌を切る時点を遅らせる理由は、自分の手牌は門前が厳しい状況で、役牌の重なりを狙うため。これはドラ3やアガリトップなどアガりたい局面。もう一つ役牌を切るのを遅らせる理由は\textbf{自分の手が悪しぎて、相手を鳴かせて横移動や局を安くて流されることを狙うため}\citep{mark2-shibori}。「絞り」と違う点は、序盤で切らなかった生牌役牌を、\textbf{7巡目あたりで放つ}。巡目が深いほど他家の手牌に役牌が重なる確率も高くなる。中盤で誰も大きな動きがない状況で、一人に鳴かせることによって、安く済むことや横移動が期待できる。たとえ立体図~\ref{rittai:shiboru-2} は生牌の \mj{67z} を7巡目までキープしてここで役牌を切り出す。

\begin{rittai}[h]
  \centering  \drawrittai{
    jicha/seat = 西,
    jicha/hand = 229m 189p 27s 122467z,
    jicha/river = 24p 5m 4s 58m,
    %
    toimen/river = 1z 1m 5z 9m 4z 7s,
    toimen/riverb = 3p,
    %
    kami/river = 1z 1m 1s 5z 4z 7p,
    kami/riverb = 2s,
    %
    shimo/river = 1z 9s 9p 5z 8m 2s,
    %
    dora = 4z,
  }
  \caption{クソ配牌をもらい、七対子や役牌を見つつ中盤に入ったら他家に鳴かせる進行になる。}
  \label{rittai:shiboru-2}
\end{rittai}

もちろん、鳴かせた相手の手牌にドラ3が秘蔵されているケースもあるが、副露の部分に見えない限りそういったケースは極端のケース。ドラ3よりドラ0やドラ1の方が多いので、鳴かせるに値する。ドラポンされていて、役牌バックが濃厚な場合や、混一色やトイトイの模様に見える場合なら絞る。


\subsubsection*{めくり合いに持ち込むための切り遅れ}

もう一つ役牌を切るのを遅らせる理由は、自分が聴牌するとき相手を鳴かせてめくり合いに持ち込むため\citep[][p.~19]{zero-teyaku}。節~\ref{subsub:yakuhai} の原則でいくと、手牌~\ref{mj:shiboru-3} の \mj{7z} は真っ先に切るべき。しかし、ツモが効かないまま生牌役牌を中盤まで抱えることも珍しくないだろう。それに、中盤で \ref{mj:shiboru-3a} のような役牌バック濃厚な仕掛けが入っている相手がいる。こうなると \mj{7z} ではなく打 \mj{3p} とする。

\tehai{34067m 33488p 234s 7z}[][shiboru-3]

\tehai{3z1m9p7s4s2s 1m2p -99'9s}[他家の河と仕掛け　ドラ \mj{9m}][shiboru-3a]

中盤なら一向聴や聴牌の可能性が高い。\mj{3p7z} のどちらを切ってもロンされてもおかしくない。問題は一向聴だった場合。相手の手牌は~\ref{mj:shiboru-3b} だったら、\mj{7z} を鳴かせない限り、打たれた \mj{58m} にロンの声をかけることができない。\mj{7z} を聴牌まで絞って、相手が自分の先にアガることを阻止し、相手を放銃の抽選に参加させることにもなる。

\tehai{6799m 234p 477z -99'9s}[ドラ \mj{9m}][shiboru-3b]

もちろん、\mj{58m} が先に入ってしまい、\mj{7z} で振り込むケースもあるので、この切り遅れは、放銃率が上がる選択とも言える。放銃率が高まるのに伴い、自分のアガリ率も高まる。


\renewcommand{\renderOption}{0}






\end{document}